%% Copyright 2007-2025 Elsevier Ltd
%%
%% This file is part of the 'Elsarticle Bundle'.
%% ---------------------------------------------
%%
%% It may be distributed under the conditions of the LaTeX Project Public
%% License, either version 1.3 of this license or (at your option) any
%% later version. The latest version of this license is in
%%    http://www.latex-project.org/lppl.txt
%% and version 1.3 or later is part of all distributions of LaTeX
%% version 2005/12/01 or later.

\documentclass[preprint,12pt]{elsarticle}

\usepackage{amssymb}
\usepackage{amsmath}
\usepackage{lineno,hyperref}
\modulolinenumbers[5]
\usepackage{graphicx}
\usepackage{booktabs}
\usepackage{multirow}
\usepackage{xcolor}
\usepackage{array}
\usepackage{tabularx}
\usepackage{algorithm}
\usepackage{algpseudocode}
\usepackage{subcaption}

\journal{Biomedical Signal Processing and Control}

\bibliographystyle{elsarticle-num}

% FiLM 5-class results (scripts/5class_extra_baselines.py, seed=42)
\newcommand{\filmfiveclassacc}{96.88}
\newcommand{\filmfiveclassci}{[95.8, 97.7]}
\newcommand{\filmfiveclassp}{0.061}

\begin{document}

\begin{frontmatter}

\title{Spectroscopy-Informed Dual-Stream Raman-Wavelet Fusion for Robust Bacterial Classification}

\author{}
\address{}
\begin{abstract}
Rapid and accurate identification of pathogenic bacteria is critical
for clinical diagnostics and antimicrobial resistance (AMR) management.
Raman spectroscopy provides label-free molecular fingerprints within
minutes, but current deep learning approaches process spectra as
one-dimensional vectors, discarding the multi-scale structural
information available in the scale-wavenumber domain and remaining
sensitive to acquisition artifacts present in portable instruments.
This work proposes a spectroscopy-informed dual-stream framework that
jointly processes raw 1D Raman spectra
and their 2D continuous wavelet transform (CWT) scalograms.
Evaluated on the publicly available Bacteria-ID benchmark across
five clinically significant species, SDR-Fusion achieves 95.84\%
test accuracy (95\% CI: [94.6, 96.8]), decisively outperforming the
spectral-only baseline ($+30.80$~pp, $p<0.001$) and providing competitive
accuracy against the strong scalogram-only ($95.20\%$) and naive
concatenation ($94.80\%$) baselines while uniquely producing
per-modality gating signals that expose modality discordance. The core
architectural contribution is a separation-aware fusion module that
gates each modality independently before cross-attention. This
separation principle prevents noise in one acquisition modality from
corrupting the other, yielding 13--18 percentage points higher accuracy
under Gaussian noise (flat 95.8\% vs.\ 77.5\% at 20\,dB SNR), and a
$+5.0$ percentage point advantage at $20\%$ training data over the
no-fusion dual-stream baseline. With 350\,K parameters and 2.3\,ms inference
latency, SDR-Fusion satisfies practical constraints for embedded Raman
inference. Integrated-gradient and layer-wise relevance propagation maps applied
to the spectral branch localize attribution weight to biochemically
established bands---amide~I/III (protein-targeting antibiotic
mechanisms), nucleic-acid stretches (fluoroquinolone and rifampicin
targets), and the pyocyanin band (\textit{Pseudomonas aeruginosa}
virulence)---providing a direct computational link between pathogen
identification and antibiotic mechanism selection. SDR-Fusion provides comparable accuracy to strong multimodal baselines while offering interpretable modality-discordance signals.
\end{abstract}
\begin{keyword}
Bacterial classification \sep
Raman spectroscopy \sep
Multimodal fusion \sep
Continuous wavelet transform \sep
Explainable AI \sep
Biomedical signal processing \sep
Antimicrobial resistance \sep
Deep learning
\end{keyword}

\end{frontmatter}

\linenumbers

% ======================================================================
\section{Introduction}
\label{sec:intro}
% ======================================================================

Rapid and accurate identification of bacterial pathogens is a central
challenge in clinical microbiology, public health, and food
safety~\cite{ho2019rapid,singh2022culture}. The global spread of
antimicrobial resistance (AMR) demands diagnostic tools that
substantially outpace culture-based methods (24--72~hours) while
remaining more accessible than whole-genome
sequencing~\cite{singh2022culture}. Raman
spectroscopy provides label-free, non-destructive molecular
fingerprints within minutes by recording inelastic scattering of a
laser source~\cite{ho2019rapid,zhou2022ramannet}. The vibrational
signatures reflect the molecular composition of key bacterial
biomolecules---proteins (amide I/III at 1600--1700 and
1200--1350~cm$^{-1}$), nucleic acids ($\sim$800--900~cm$^{-1}$),
and lipids (1000--1200~cm$^{-1}$)---forming a chemically
interpretable fingerprint (Fig.~\ref{fig:single_spectrum}) that
distinguishes bacterial species
(Fig.~\ref{fig:class_spectra})~\cite{lussier2020deep,singh2022culture}.

\begin{figure*}[!t]
\centering
\begin{subfigure}{0.48\textwidth}
  \centering
  \includegraphics[width=\linewidth]{figures/fig1_single_spectrum.png}
  \caption{Representative Raman spectrum}
  \label{fig:single_spectrum}
\end{subfigure}\hfill
\begin{subfigure}{0.48\textwidth}
  \centering
  \includegraphics[width=\linewidth]{figures/fig2_class_spectra.png}
  \caption{Mean Raman spectra by class}
  \label{fig:class_spectra}
\end{subfigure}
\caption{(a) Representative Raman spectrum with biochemical bands.
(b) Mean Raman spectra for five species, showing spectral overlaps
and distinctive signatures (e.g., pyocyanin).}
\label{fig:combined_spectra}
\end{figure*}

These molecular fingerprints are directly relevant to therapeutic
strategy: amide-band alterations indicate protein-targeting antibiotic
mechanisms (e.g., $\beta$-lactams disrupting cell-wall peptidoglycan),
nucleic-acid band shifts point toward nucleic-acid-targeting agents
(fluoroquinolones, rifampicin), and the pyocyanin band uniquely
identifies \textit{P.~aeruginosa}, for which specific anti-virulence
strategies exist~\cite{singh2022culture}. An AI system that can both
\emph{identify} a pathogen and \emph{localize} the spectral bands
driving its predictions can provide attribution maps that offer
molecular-level cues for downstream therapeutic analysis.

Furthermore, the global AMR crisis increasingly demands diagnostic
intelligence at the network edge: rural clinics, intensive care units,
and field hospitals cannot rely on centralized laboratory
infrastructure~\cite{maruthamuthu2020rapid}. Point-of-care Raman
spectrometers are now commercially available at sub-laboratory price
points accessible to mid-tier clinical settings~\cite{ho2019rapid},
but lack embedded
AI classifiers that are simultaneously accurate, noise-robust, and
lightweight enough for FPGA or microcontroller deployment. We therefore investigate whether robustness-aware multimodal fusion can improve bacterial Raman classification under realistic acquisition variability.

\subsection{Gap in Existing Approaches}
\label{ssec:gap}

Conventional deep-learning approaches for Raman classification treat spectra as 1D vectors~\cite{ho2019rapid,zhou2022ramannet,
balytskyi2024open}; while computationally efficient, this constraint limits feature extraction to the amplitude domain. Liu et~al.\ converted Raman spectra to
2D wavelet-packet Gramian angular field images for CNN
classification~\cite{liu2024wpga}, demonstrating the value of
scale-domain representations; however, this approach \emph{replaces}
the 1D spectral branch, discarding precise peak-amplitude information.
Chia et~al.\ used wavelet knockoff features for interpretable
classification~\cite{chia2021interpretable}, and while both approaches
demonstrate scale-domain value, neither fuses wavelet features with the
raw 1D signal in a jointly trained, adaptive dual-stream architecture.
To the best of our knowledge, prior Raman classification studies have not jointly trained raw spectral and CWT representations using adaptive per-sample fusion.

Multimodal fusion architectures from medical imaging (CT--MRI,
histology--genomics) are designed for \emph{heterogeneous} data
sources whose latent spaces are statistically
independent~\cite{holzinger2022information}. In Raman--CWT fusion,
by contrast, every scalogram feature is \emph{deterministically
derived} from the corresponding spectral sample, making the two
representations physically aligned. 

Current AI-assisted spectroscopic models rarely map learned
attributions directly back to domain-specific pharmacological
mechanisms or actionable clinical targets. Similarly, most existing
high-accuracy Raman classifiers are either too computationally heavy
for edge deployment or lack the noise robustness required to handle
lower-quality spectra from portable IoT instruments. 

\subsection{Contributions}
\label{ssec:contributions}

SDR-Fusion addresses these gaps with three specific, evidence-backed
contributions:

\begin{enumerate}

\item \textbf{SeparationCross fusion block.}
      Per-modality learnable sigmoid gating (conditioned
      independently on each branch before any cross-modal interaction)
      followed by asymmetric cross-attention (spectral query,
      scalogram key/value). While gated multimodal fusion has been
      explored for heterogeneous data~\cite{arevalo2017gated}, the
      true distinguishing property of SeparationCross is the
      \emph{separation principle}---that each gate is conditioned
      only on its own modality before any cross-modal interaction.
      This separation prevents noise in one modality from contaminating
      the gating of the other.

\item \textbf{Spectral attribution.}
      Integrated-gradient and LRP maps applied to the 1D spectral
      branch produce class-wise wavenumber importance profiles that
      identify the specific molecular bands driving each prediction.
      These maps provide domain-specific interpretation mapping bands
      to antibiotic mechanisms (e.g., amide I/III, nucleic-acid
      stretches, pyocyanin), providing interpretable evidence that the
      classification is driven by physically meaningful features
      rather than spurious correlations.

\item \textbf{Robustness and data-scarcity resilience.}
      Under Gaussian noise (SNR 20--40~dB), baseline drift, band
      occlusion, and data scarcity, SDR-Fusion maintains consistent
      margins over single-modality baselines. At 20\% training
      data---a realistic constraint in resource-limited clinical
      settings---SDR-Fusion outperforms Raman-only and no-fusion
      baselines, demonstrating the regularizing effect of structurally
      aligned fusion.

\end{enumerate}

% ======================================================================
\section{Related Work}
\label{sec:related}

% ======================================================================
\subsection{Raman Spectroscopy and Deep Learning for Bacterial
Identification}
\label{ssec:raman_dl}

Early Raman-based bacterial identification relied on classical machine
learning---PCA-LDA, SVM, and random forests---achieving 70--85\% on
small datasets~\cite{ho2019rapid,zhou2022ramannet}. Deep learning
substantially improved performance: Ho et~al.\ achieved 82.2\% with
ResNet-18~\cite{ho2019rapid}; Zhou et~al.\ reached 84.7\% with
RamanNet, a 30\,k-parameter 1D CNN~\cite{zhou2022ramannet}; Singh
et~al.\ applied deep transfer learning to clinical
samples~\cite{singh2022culture}; Balytskyi et~al.\
extended to open-world identification with Objectosphere
regularization~\cite{balytskyi2024open}. Liu et~al.\ converted spectra
to 2D wavelet-packet images, achieving 90.6\%~\cite{liu2024wpga};
however, this approach replaces rather than supplements the 1D branch,
losing amplitude-domain peak precision. More recently, transformer-based
architectures have been applied to vibrational spectra~\cite{yan2023spectral},~\cite{lussier2020deep}. Despite these advances, to the best of our knowledge, no prior
work jointly trains a 1D spectral branch and a 2D CWT branch with
adaptive, per-sample modality weighting.

% ======================================================================
\subsection{Wavelet and Scale-Domain Representations for Spectroscopic
Analysis}
\label{ssec:wavelet}

Wavelet transforms underpin multi-resolution analysis across biomedical
signal processing, including ECG denoising, EEG artifact
removal, and spectroscopic
preprocessing~\cite{lussier2020deep,maruthamuthu2020rapid}. Wavelet-derived spectral representations have also shown strong
performance for bacterial Raman classification and biochemical feature
localization~\cite{liu2024wpga}. For Raman
data specifically, Chia et~al.\ used wavelet knockoff features for
statistically controlled, interpretable
classification~\cite{chia2021interpretable}, and Liu et~al.\ used
wavelet-packet images as CNN input~\cite{liu2024wpga}. Both studies
demonstrate the value of scale-domain representations, but neither
retains the raw 1D spectral branch alongside the wavelet-derived
features in a jointly optimised architecture. The continuous wavelet
transform (CWT) with a real Morlet wavelet offers logarithmic
frequency resolution that is particularly well matched to the broad
baseline envelopes and sharp localised peaks that co-exist in bacterial
Raman fingerprints~\cite{chia2021interpretable,liu2024wpga}; this
property motivates its use in SDR-Fusion over fixed-resolution
alternatives such as the short-time Fourier transform.

% ======================================================================
\subsection{Multimodal and Multi-Stream Fusion in Biomedical AI}
\label{ssec:fusion}

Holzinger et~al.\ establish information fusion as an integrative enabler
for robust, explainable, and trustworthy medical
AI~\cite{holzinger2022information}. Within this paradigm, gated
multimodal units~\cite{arevalo2017gated} introduced learnable sigmoid
gates for heterogeneous modality fusion, and feature-wise linear
modulation (FiLM)~\cite{perez2018film} demonstrated expressive affine
conditioning across modalities. Attention-based fusion has been explored
in multimodal medical imaging---for example, combining CT and MRI
streams via cross-attention for tumour
segmentation~\cite{holzinger2022information}---and in audio-visual
learning via asymmetric cross-modal
attention~\cite{vaswani2017attention}. However, these generic fusion
modules assume statistically independent latent spaces. In Raman--CWT
fusion, every scalogram feature is deterministically derived from the
same spectral sample, making the two representations physically aligned;
this structural property demands a fusion design where each branch's
gate is conditioned only on its own signal quality before any
cross-modal interaction, which is the distinguishing principle of
SeparationCross. Ma et~al.\ address adversarial robustness in medical
image analysis~\cite{ma2021understanding}; unlike their synthetic-attack
focus, SDR-Fusion targets the physical noise and baseline drift inherent
in point-of-care Raman acquisition. Gated Multimodal Units (GMU) introduced cross-conditioned gating for
multimodal fusion through shared modality interaction
mechanisms~\cite{arevalo2017gated}.

% ======================================================================% ======================================================================
\subsection{Interpretable and Deployable Raman AI}
\label{ssec:interp_edge}
% ======================================================================

Integrated gradients~\cite{sundararajan2017axiomatic} and layer-wise
relevance propagation (LRP)~\cite{bach2015pixel} are established
methods for attributing neural-network predictions to input features.
In Raman spectroscopy, attribution maps have been used to identify
chemically meaningful bands and validate that attend to
biologically relevant spectral regions rather than acquisition
artefacts~\cite{lussier2020deep,singh2022culture,chia2021interpretable}.
SDR-Fusion extends this paradigm within a dual-stream architecture by
mapping class-specific predictions to interpretable Raman bands. Interpretability-driven Raman learning has additionally been explored
through knockoff-wavelet attribution frameworks on the same
Bacteria-ID benchmark~\cite{chia2021interpretable}.

Beyond interpretability, Raman-AI systems must also satisfy practical
deployment constraints including low latency, compact parameterization,
and robustness to acquisition variability. Model compression
techniques such as quantization and structured
pruning~\cite{han2016deep,jacob2018quantization} have enabled compact
deep-learning models for embedded and edge-AI deployment, motivating
lightweight multimodal architectures for portable Raman diagnostics.
%  ======================================================================
\section{Methodology}
\label{sec:methods}
% ======================================================================

\subsection{Notations and Problem Formulation}
\label{ssec:notation}

Table~\ref{tab:notations} summarizes the mathematical symbols used
throughout. Let $\mathcal{D}=\{(\mathbf{x}_i^{(s)},\mathbf{x}_i^{(w)},
y_i)\}_{i=1}^{N}$, where $\mathbf{x}_i^{(s)}\in\mathbb{R}^{L}$
($L=1000$) is the preprocessed Raman spectrum spanning
400--1800~cm$^{-1}$, $\mathbf{x}_i^{(w)}\in\mathbb{R}^{3\times224\times224}$
is the corresponding CWT scalogram, and $y_i\in\{0,\ldots,C{-}1\}$
is the bacterial class label ($C=5$). The objective is to learn
$f:(\mathbf{x}^{(s)},\mathbf{x}^{(w)})\rightarrow y$ that maximizes
accuracy on unseen samples while producing attribution maps that are
interpretable in terms of established Raman band assignments.

\begin{table}[!t]
\centering
\caption{Mathematical notations used in the methodology.}
\label{tab:notations}
\renewcommand{\arraystretch}{1.2}
\begin{tabular}{lp{4.2cm} | lp{4.2cm}}
  \toprule
  \textbf{Symbol} & \textbf{Description} &
  \textbf{Symbol} & \textbf{Description} \\
  \midrule
  $N$ & Number of bacterial samples &
  $\alpha,\,\beta$ & Learnable sigmoid gates $\in(0,1)^{128}$ \\
  $L$ & Spectrum length ($L=1000$) &
  $\mathbf{W}_\alpha,\mathbf{W}_\beta$ &
      Gating weight matrices $\in\mathbb{R}^{128\times128}$ \\
  $\mathbf{x}_i^{(s)}$ & Preprocessed 1D Raman spectrum &
  $\mathbf{b}_\alpha,\mathbf{b}_\beta$ &
      Gating bias vectors $\in\mathbb{R}^{128}$ \\
  $\mathbf{x}_i^{(w)}$ & CWT scalogram ($3\times224\times224$) &
  $\mathbf{Q},\mathbf{K},\mathbf{V}$ & Query, Key, Value for attention \\
  $y_i$ & Ground-truth label $\in\{0,\ldots,C{-}1\}$ &
  $d_k$ & Attention scaling dim ($d_k=128$) \\
  $C$ & Number of bacterial classes ($C=5$) &
  $\mathbf{z}_\mathrm{fused}$ & 256-dim fused representation \\
  $\mathbf{z}^{(s)},\mathbf{z}^{(w)}$ &
      128-dim spectral/scalogram embeddings &
  $\mathcal{L}_\mathrm{CE}$ & Cross-entropy training loss \\
  \bottomrule
\end{tabular}
\end{table}

\subsection{Data Preprocessing}
\label{ssec:preprocessing}

All preprocessing is applied identically to all model variants and was
fixed before any model evaluation. Table~\ref{tab:preprocparams}
provides a complete parameter listing for reproducibility.

\subsubsection{1D Raman Spectra}
\label{sssec:prep1d}

Three sequential steps clean the raw spectrum:

\textbf{(1) Baseline correction.} Asymmetric least-squares (ALS)
smoothing~\cite{lussier2020deep} estimates and subtracts the
fluorescence background.

\textbf{(2) Normalization.} Min--max scaling maps each corrected
spectrum to $[0,1]$.

\textbf{(3) Denoising.} Savitzky--Golay filtering (second-order
polynomial, window 11) suppresses high-frequency noise while
preserving peak shapes.

\subsubsection{CWT Scale--Wavenumber Scalogram Generation}
\label{sssec:cwt}

A real Morlet CWT (implemented via PyWavelets v1.4.1~\cite{lee2019pywavelets},
wavelet identifier \texttt{morl}) is applied to each preprocessed spectrum:
\begin{equation}
  W(a,b) = \frac{1}{\sqrt{a}}
  \int_{-\infty}^{\infty} x(t)\,\psi^*\!\left(\frac{t-b}{a}\right)
  \mathrm{d}t,
  \label{eq:cwt}
\end{equation}
where $a>0$ is the scale parameter and $b$ is the translation
(wavenumber position, cm$^{-1}$). Because the Raman spectrum's
horizontal axis is wavenumber rather than time, the output is a
\emph{scale--wavenumber} map. The power scalogram $|W(a,b)|^2$ is
computed over $N_s{=}256$ log-spaced scales
($a_i = 2^{(i-1)\cdot\log_2(256)/(255)}$, $i=1,\ldots,256$)
via \texttt{pywt.cwt(spectrum, scales, 'morl')}.
The resulting $256{\times}1000$ power map is resized to
$224{\times}224$ pixels by bilinear interpolation
(\texttt{cv2.resize}, \texttt{INTER\_LINEAR}), normalised to $[0,1]$,
quantised to 8-bit, and mapped to a three-channel Viridis RGB image
via \texttt{cv2.applyColorMap(\ldots, cv2.COLORMAP\_VIRIDIS)}
(Fig.~\ref{fig:wavelet}).

\begin{figure}[!t]
\centering
\includegraphics[width=0.48\textwidth]{figures/fig3_wavelet_example.png}
\caption{CWT converts a 1D Raman spectrum into a 2D scalogram,
capturing complementary scale-wavenumber features.}
\label{fig:wavelet}
\end{figure}

\textbf{Physical and pharmacological justification.}
The fingerprint region 400--1800~cm$^{-1}$ encompasses all clinically
established Raman band assignments for bacterial biomolecules:
amide~I/III (protein secondary structure, relevant to $\beta$-lactam
and glycopeptide mechanisms), nucleic-acid C--O stretches
(fluoroquinolone and rifampicin targets), phenylalanine
($\sim$1003~cm$^{-1}$), and lipid C--C/C--H stretches (polymyxin and
daptomycin membrane-disruption targets)~\cite{singh2022culture}.
Confining the CWT to this region not only suppresses Rayleigh-line
artefacts but ensures that every scale--wavenumber feature has a
chemically interpretable correspondent in the Raman assignment
literature---directly supporting biochemically interpretable
attribution.

\subsection{SDR-Fusion Architecture}
\label{ssec:architecture}

SDR-Fusion comprises four components: (1) a 1D CNN spectral branch,
(2) a 2D CNN scalogram branch, (3) the SeparationCross fusion module,
and (4) a classifier head. Fig.~\ref{fig:arch_combined} shows the
complete dual-stream architecture including branch-level details.

\begin{figure*}[!t]
\centering
\includegraphics[width=\textwidth,keepaspectratio]{Raman Arch.jpeg}
\caption{Consolidated SDR-Fusion architecture for Raman--CWT fusion. A dual-stream model processes the 1D Raman spectrum and 2D CWT scalogram via dedicated CNN branches, followed by SeparationCross fusion with independent gating ($\alpha$, $\beta$), cross-attention, and residual scaling ($\gamma$). The fused representation is passed to a classifier head for prediction. Total parameters: $\sim$350K.}
\label{fig:arch_combined}
\end{figure*}

\subsubsection{1D Spectral Branch}
\label{sssec:branch1d}

The 1D spectral branch is designed to extract hierarchical
morphological features directly from the preprocessed Raman spectrum
using a stride-2 stem followed by a residual block and a projection
layer. The stride-2 stem halves the sequence length early, providing
a compact multi-scale representation:
\begin{align}
  \mathbf{h}_1^{(s)} &= \mathrm{ReLU}\!\left(\mathrm{BN}\!\left(
      \mathrm{Conv1D}(\mathbf{x}^{(s)};64, k{=}7, s{=}2)\right)\right).
\end{align}
A residual block then builds a 128-channel representation with a
skip connection that preserves gradient flow to early layers:
\begin{align}
  \mathbf{r}^{(s)} &= \mathrm{BN}\!\left(
      \mathrm{Conv1D}(\mathbf{h}_1^{(s)};128, k{=}5)\right), \\
  \mathbf{h}_2^{(s)} &= \mathrm{ReLU}\!\left(
      \mathrm{BN}\!\left(\mathrm{Conv1D}(\mathrm{ReLU}(\mathbf{r}^{(s)});128, k{=}3)\right)
      + \mathbf{r}^{(s)}\right).
\end{align}
A final projection layer maps to the embedding dimension:
\begin{align}
  \mathbf{h}_3^{(s)} &= \mathrm{ReLU}\!\left(\mathrm{BN}\!\left(
      \mathrm{Conv1D}(\mathbf{h}_2^{(s)};128, k{=}3)\right)\right).
\end{align}
Global adaptive average pooling then aggregates the spatial dimension,
yielding a compact, fixed-length spectral embedding
$\mathbf{z}^{(s)}\in\mathbb{R}^{128}$ that remains robust to minor
peak shifts while retaining critical amplitude-domain information.
The residual connection ensures that gradients reach the stem directly,
which is essential for learning from the 1000-point input.
\subsubsection{Spectral Tokenization and Feature Embedding}
\textit{Note: the tokenization described here is performed inside the
SeparationCross fusion module (Section~\ref{sssec:fusion}) as the
preparatory step that maps CNN branch outputs into aligned token
sequences for cross-attention. It does not replace or interpose the
CNN branches described above; the 1D and 2D CNN encoders run to
completion first, producing embeddings $\mathbf{z}^{(s)}$ and
$\mathbf{z}^{(w)}$, and tokenization is then applied to the
intermediate feature maps before the attention operation.}

Inside the SeparationCross module, each CNN branch output is first
modality-gated and then adaptively pooled into a fixed-length token
sequence. Both branches share embedding dimension $D{=}128$.

For the spectral branch, the per-channel sigmoid gate
$\alpha \in [0,1]^{D}$ (derived from the pooled spectral embedding;
see Section~\ref{sssec:fusion}) scales the pre-pooling 1D CNN feature
map $\mathbf{h}^{(s)} \in \mathbb{R}^{D \times L'}$.
The gated map $\mathbf{h}_\alpha^{(s)} = \alpha \odot \mathbf{h}^{(s)}$
is adaptively pooled to $S$ wavenumber-aligned tokens:
\begin{equation}
    \mathbf{T}^{(s)} =
      \bigl(\mathrm{AdaptiveAvgPool1d}_{S}\!\bigl(\mathbf{h}_\alpha^{(s)}
      \bigr)\bigr)^{\!\top} \in \mathbb{R}^{S \times D},
\end{equation}
where $S{=}25$ spectral tokens are selected by grid search on the
validation set (Section~\ref{sssec:fusion}).

Concurrently, the per-channel gate $\beta \in [0,1]^{D}$ scales the
2D CNN feature map $\mathbf{h}^{(w)} \in \mathbb{R}^{D \times H' \times W'}$.
The gated map $\mathbf{h}_\beta^{(w)} = \beta \odot \mathbf{h}^{(w)}$
is spatially flattened into a token sequence:
\begin{equation}
    \mathbf{T}^{(w)} =
      \bigl(\mathrm{Flatten}_{H'W'}\!\bigl(\mathbf{h}_\beta^{(w)}
      \bigr)\bigr)^{\!\top} \in \mathbb{R}^{H'W' \times D}.
\end{equation}
Token index $i$ in $\mathbf{T}^{(s)}$ aligns with the same relative
spectral band window as index $i$ in $\mathbf{T}^{(w)}$, preparing
both streams for cross-modal attention in the SeparationCross module.
\subsubsection{2D Scalogram Branch}
\label{sssec:branch2d}

In parallel, the 2D scalogram branch processes the continuous wavelet
transform output to capture multi-scale spatial patterns, which encode
peak widths and noise-robust structural invariants. This branch
utilizes three stacked 2D convolutional blocks, each featuring a
standard $3\times3$ spatial kernel, batch normalization, and a
$\mathrm{ReLU}$ activation. To systematically reduce the spatial
resolution while extracting dominant scale-domain features, a maximum
pooling operation ($\mathrm{MaxPool2D}$) is applied after each
activation:
\begin{align}
  \mathbf{h}_1^{(w)} &= \mathrm{MaxPool2D}\!\left(\mathrm{ReLU}\!\left(
      \mathrm{BN}\!\left(\mathrm{Conv2D}(\mathbf{x}^{(w)};32,
      3{\times}3)\right)\right)\right), \\
  \mathbf{h}_2^{(w)} &= \mathrm{MaxPool2D}\!\left(\mathrm{ReLU}\!\left(
      \mathrm{BN}\!\left(\mathrm{Conv2D}(\mathbf{h}_1^{(w)};64,
      3{\times}3)\right)\right)\right), \\
  \mathbf{h}_3^{(w)} &= \mathrm{MaxPool2D}\!\left(\mathrm{ReLU}\!\left(
      \mathrm{BN}\!\left(\mathrm{Conv2D}(\mathbf{h}_2^{(w)};128,
      3{\times}3)\right)\right)\right).
\end{align}
Adaptive average pooling then collapses the remaining spatial
dimensions to produce the scalogram embedding
$\mathbf{z}^{(w)}\in\mathbb{R}^{128}$. This ensures the output
dimensionality structurally matches the 1D spectral branch, preparing
the two distinct representations for modality-aware gating in the
subsequent SeparationCross module.

\subsection{SeparationCross Fusion Module}
\label{sssec:fusion}

Before cross-attention, the gated feature maps are adaptively pooled into fixed-length token sequences to enable aligned spectral–scalogram interaction. The SeparationCross module (Fig.~\ref{fig:arch_combined}d) implements
two complementary mechanisms designed for physically aligned 1D--2D
spectroscopic fusion, drawing inspiration from gated
mixture-of-experts~\cite{shazeer2017outrageously}. Optimal fusion should be \emph{asymmetric} and \emph{per-sample quality-aware}: the 1D
branch provides sharp peak localization, the CWT branch provides
noise-robust multi-scale context, and a sample-wise gate should route
trust toward whichever branch is more reliable. Existing fusion modules
assume statistical independence between modality embeddings; for
Raman--CWT fusion, asymmetric, per-sample quality-aware
gating---absent from all generic modules---is necessary. If a standard
fusion mechanism mixes these embeddings indiscriminately, noise in a
single acquisition corrupts the entire fused embedding. SDR-Fusion avoids
this by designing a module where each branch's gate is conditioned
solely on its own signal quality.
\noindent\textbf{Step 1 --- Separation gating.}
\begin{align}
  \alpha &= \sigma\!\left(\mathbf{W}_\alpha\mathbf{z}^{(s)}
           +\mathbf{b}_\alpha\right), \label{eq:alpha}\\
  \beta  &= \sigma\!\left(\mathbf{W}_\beta\mathbf{z}^{(w)}
           +\mathbf{b}_\beta\right).  \label{eq:beta}
\end{align}
Gated embeddings: $\mathbf{z}_\alpha^{(s)}=\alpha\odot\mathbf{z}^{(s)}$,
$\mathbf{z}_\beta^{(w)}=\beta\odot\mathbf{z}^{(w)}$. Each gate is
conditioned \emph{only} on its own modality, ensuring that spectral
signal quality controls only the spectral branch and scalogram quality
controls only the scalogram branch.

\noindent\textbf{Step 2 --- Spatial cross-attention.}
The gated feature maps are tokenized into spatial sequences. The
gated spectral feature map
$\mathbf{h}_\alpha^{(s)} = \alpha \odot \mathbf{h}^{(s)} \in \mathbb{R}^{D \times L}$
(where $\mathbf{h}^{(s)}$ is the pre-pooling 1D branch output)
is pooled to $S{=}25$ tokens via adaptive average pooling and
transposed to yield $\mathbf{T}^{(s)} \in \mathbb{R}^{S \times D}$.
The value $S{=}25$ was selected by grid-searching $S \in \{10, 15,
20, 25, 30, 50\}$ on the validation set; $S{=}25$ achieved the
highest validation accuracy and aligns with the approximate width of
individual biochemical peaks at 400--1800~cm$^{-1}$ spectral
resolution. The gated scalogram feature map
$\mathbf{h}_\beta^{(w)} = \beta \odot \mathbf{h}^{(w)} \in \mathbb{R}^{D \times H' \times W'}$
is flattened to $\mathbf{T}^{(w)} \in \mathbb{R}^{H'W' \times D}$.
Cross-attention~\cite{vaswani2017attention},~\cite{tsai2019multimodal} is computed with the
spectral tokens as queries and scalogram tokens as keys/values:
\begin{equation}
  \mathbf{Q} = \mathbf{T}^{(s)}\mathbf{W}_Q, \quad
  \mathbf{K} = \mathbf{T}^{(w)}\mathbf{W}_K, \quad
  \mathbf{V} = \mathbf{T}^{(w)}\mathbf{W}_V,
  \label{eq:qkv}
\end{equation}
where $\mathbf{W}_Q, \mathbf{W}_K, \mathbf{W}_V \in \mathbb{R}^{D
\times D}$ and $D{=}128$. The attended output is:
\begin{equation}
  \mathbf{Z}_\mathrm{attn} =
  \mathrm{softmax}\!\left(\frac{\mathbf{Q}\mathbf{K}^\top}{\sqrt{D}}
  \right)\mathbf{V} \in \mathbb{R}^{S \times D},
  \label{eq:attn}
\end{equation}
which is mean-pooled over the $S$ tokens to produce
$\mathbf{z}_\mathrm{attn} \in \mathbb{R}^{D}$. A learnable residual
scalar $\gamma$ (initialized to zero) controls the injection of
spectral evidence:
$\mathbf{z}_\mathrm{combined} = \mathbf{z}_\alpha^{(s)} + \gamma
\cdot \mathbf{z}_\mathrm{attn}$. The fused vector:
\begin{equation}
  \mathbf{z}_\mathrm{fused} =
  \left[\mathbf{z}_\mathrm{combined};\,\mathbf{z}_\beta^{(w)}\right]
  \in\mathbb{R}^{2D}.
  \label{eq:fused}
\end{equation}

\subsection{Classifier Head and Training Procedure}
\label{ssec:head_and_training}

The final classification is performed by a dedicated head that maps
the fused multimodal embedding to class probabilities. The
256-dimensional fused vector is passed through a fully connected
layer, regularized via dropout, and projected to the final class
logits using a softmax activation:
\begin{align}
  \mathbf{h}_\mathrm{cls}  &= \mathrm{ReLU}\!\left(
      \mathrm{FC}(\mathbf{z}_\mathrm{fused};256\!\to\!128)\right), \\
  \mathbf{h}_\mathrm{drop} &= \mathrm{Dropout}(
      \mathbf{h}_\mathrm{cls},p{=}0.3), \\
  \hat{\mathbf{y}} &= \mathrm{Softmax}\!\left(
      \mathrm{FC}(\mathbf{h}_\mathrm{drop};128\!\to\!C)\right).
\end{align}
The entire architecture is optimized end-to-end using the standard
categorical cross-entropy loss function over the $N$ samples and $C$
classes:
\begin{equation}
  \mathcal{L}_\mathrm{CE} =
  -\frac{1}{N}\sum_{i=1}^{N}\sum_{c=0}^{C-1}
  y_{i,c}\log\hat{y}_{i,c}.
\end{equation}
Table~\ref{tab:hyperparams} lists the optimization settings and
hyperparameters, which were fixed prior to any model evaluation. The
comprehensive end-to-end training process and the subsequent
low-latency edge-inference procedure are detailed sequentially in
Algorithms~\ref{alg:training} and \ref{alg:inference}.

\begin{table}[!t]
\centering
\begin{minipage}[t]{0.48\textwidth}
\centering
\caption{Preprocessing pipeline parameters.}
\label{tab:preprocparams}
\renewcommand{\arraystretch}{1.2}
\resizebox{\textwidth}{!}{
\begin{tabular}{lll}
  \toprule
  Step & Parameter & Value \\
  \midrule
  ALS baseline correction & Smoothness $\lambda$ & $10^{5}$ \\
                          & Asymmetry $p$        & $0.01$ \\
                          & Iterations           & 10 \\
  Savitzky--Golay denoise & Window size          & 11 \\
                          & Polynomial order     & 2 \\
  Normalization           & Method                & Min--max to $[0,1]$ \\
  CWT (real Morlet)       & Centre freq.\ $\omega_0$ & 6~rad \\
                          & Scale range          & 1--256 (256 scales) \\
                          & Output size          & $224\times224$~px \\
                          & Color map            & Viridis (3-ch.\ RGB) \\
  \bottomrule
\end{tabular}
}
\end{minipage}\hfill
\begin{minipage}[t]{0.48\textwidth}
\centering
\caption{Training hyperparameters.}
\label{tab:hyperparams}
\renewcommand{\arraystretch}{1.2}
\resizebox{\textwidth}{!}{
\begin{tabular}{ll}
  \toprule
  Hyperparameter & Value \\
  \midrule
  Optimiser             & AdamW \\
  Initial learning rate & $1\times10^{-3}$ \\
  Weight decay          & $1\times10^{-4}$ \\
  Batch size            & 32 \\
  Max.\ epochs          & 30 \\
  LR schedule           & Cosine annealing ($T_{\max}=30$) \\
  Dropout               & 0.3 \\
  Random seed           & 42 \\
  Mixed precision       & FP16 (\texttt{torch.autocast}) \\
  Framework             & PyTorch~2.0 \\
  Hardware              & Apple Silicon (M-series, MPS backend) \\
  \bottomrule
\end{tabular}
}
\end{minipage}
\end{table}

\begin{algorithm}[!t]
\caption{End-to-end training of SDR-Fusion}
\label{alg:training}
\begin{algorithmic}[1]
  \State \textbf{Input:}
         $\mathcal{D}=\{(\mathbf{x}_i^{(s)},\mathbf{x}_i^{(w)},
         y_i)\}_{i=1}^N$;
         hyperparameters from Table~\ref{tab:hyperparams}
  \State \textbf{Output:} Trained parameters $\Theta^*$
  \State Initialize $\Theta=\{\theta_s,\theta_w,\theta_f,\theta_c\}$;
         seed $\leftarrow 42$
  \For{epoch $= 1$ \textbf{to} $30$}
    \State Shuffle $\mathcal{D}$
    \For{each mini-batch
         $\{(\mathbf{x}_b^{(s)},\mathbf{x}_b^{(w)},y_b)\}$}
      \State $\mathbf{x}_b^{(s)}\leftarrow
              \mathrm{ALS}(\cdot){\to}\mathrm{MinMax}(\cdot)
              {\to}\mathrm{SG}(\cdot)$
      \State $\mathbf{x}_b^{(w)}\leftarrow
              \mathrm{CWT}(\mathbf{x}_b^{(s)}){\to}
              \mathrm{Viridis\_RGB}(3{\times}224{\times}224)$
      \State $\mathbf{z}^{(s)}\leftarrow
              f_{1D}(\mathbf{x}_b^{(s)};\theta_s)$;\;
              $\mathbf{z}^{(w)}\leftarrow
              f_{2D}(\mathbf{x}_b^{(w)};\theta_w)$
      \State $\alpha\leftarrow
              \sigma(\mathbf{W}_\alpha\mathbf{z}^{(s)}{+}
              \mathbf{b}_\alpha)$;\;
              $\beta\leftarrow
              \sigma(\mathbf{W}_\beta\mathbf{z}^{(w)}{+}
              \mathbf{b}_\beta)$
              \hfill{\small(Eqs.~\ref{eq:alpha}--\ref{eq:beta})}
      \State $\mathbf{z}_\alpha^{(s)}\leftarrow
              \alpha\odot\mathbf{z}^{(s)}$;\;
              $\mathbf{z}_\beta^{(w)}\leftarrow
              \beta\odot\mathbf{z}^{(w)}$
      \State $\mathbf{Q}{\leftarrow}\mathbf{T}^{(s)}$;\;
              $\mathbf{K}{\leftarrow}\mathbf{T}^{(w)}$;\;
              $\mathbf{V}{\leftarrow}\mathbf{T}^{(w)}$
              \hfill{\small(tokenized via
              Sec~\ref{sssec:fusion})}
      \State $\mathbf{z}_\mathrm{attn}\leftarrow
              \mathrm{mean\_pool}(
              \mathrm{softmax}(\mathbf{Q}\mathbf{K}^\top/
              \sqrt{d_k})\mathbf{V})$
      \State $\mathbf{z}_\mathrm{combined}\leftarrow
              \mathbf{z}_\alpha^{(s)}{+}\gamma\mathbf{z}_\mathrm{attn}$
      \State $\mathbf{z}_\mathrm{fused}\leftarrow
              [\mathbf{z}_\mathrm{combined};\,
              \mathbf{z}_\beta^{(w)}]$
      \State $\hat{y}_b\leftarrow
              h_\mathrm{cls}(\mathbf{z}_\mathrm{fused};\theta_c)$
      \State $\mathcal{L}\leftarrow
              \mathcal{L}_\mathrm{CE}(\hat{y}_b,y_b)$;\;
              $\Theta\leftarrow\Theta{-}\eta_t\nabla_\Theta\mathcal{L}$
    \EndFor
  \EndFor
  \State \Return $\Theta^*$
\end{algorithmic}
\end{algorithm}

\begin{algorithm}[!t]
\caption{Inference on a new Raman spectrum (2.3~ms end-to-end)}
\label{alg:inference}
\begin{algorithmic}[1]
  \State \textbf{Input:} Raw
         $\mathbf{x}_\mathrm{new}^{(s)}\in\mathbb{R}^{1000}$
         (from any portable/IoT Raman device)
  \State \textbf{Output:} Predicted species $\hat{y}_\mathrm{new}$ +
         attribution vector $\mathbf{a}\in\mathbb{R}^{1000}$
  \State Preprocess: ALS $\to$ MinMax $\to$ SG
  \State CWT $\to$ Viridis RGB: $\mathbf{x}_\mathrm{new}^{(w)}$
  \State $\mathbf{z}^{(s)},\mathbf{z}^{(w)}\leftarrow
         f_{1D}(\cdot),f_{2D}(\cdot)$
  \State SeparationCross fusion $\to\mathbf{z}_\mathrm{fused}$
         \hfill{\small(Eqs.~\ref{eq:alpha}--\ref{eq:fused})}
  \State $\hat{y}_\mathrm{new}\leftarrow
         \arg\max_c h_\mathrm{cls}(\mathbf{z}_\mathrm{fused})[c]$
  \State $\mathbf{a}\leftarrow
         \mathrm{IntegratedGradients}(f_{1D},
         \mathbf{x}_\mathrm{new}^{(s)},\hat{y}_\mathrm{new})$
         \hfill{\small(biochemically interpretable attribution)}
  \State \Return $\hat{y}_\mathrm{new},\;\mathbf{a}$
\end{algorithmic}
\end{algorithm}


% ======================================================================
\section{Experiments and Results}
\label{sec:experiments}
% ======================================================================

\subsection{Dataset and Species Selection}
\label{ssec:dataset}
The experimental validation of the SDR-Fusion architecture is conducted on the \textit{Bacteria-ID} benchmark dataset~\cite{ho2019rapid}, which contains high-resolution single-cell Raman spectra across five clinically significant bacterial species. To ensure rigorous evaluation and prevent class imbalances from skewing baseline comparisons, the dataset partition is balanced to comprise exactly 2,500 independent spectral profiles per target category ($N = 12,500$ total spectra). The clinical profiles evaluated within this multi-class architecture include:
\begin{itemize}
    \item \textit{Escherichia coli} ($n=2,500$): A predominant Gram-negative pathogen frequently isolated as the primary cause of clinical bacteraemia and urinary tract infections (UTIs).
    \item \textit{Klebsiella pneumoniae} ($n=2,500$): An encapsulated Gram-negative rod associated with severe nosocomial pneumonia and multi-drug resistant hospital outbreaks.
    \item \textit{Pseudomonas aeruginosa} ($n=2,500$): An opportunistic Gram-negative pathogen known for its robust biofilm-forming capabilities on medical instrumentation.
    \item \textit{Staphylococcus aureus} ($n=2,500$): A critical Gram-positive species serving as a primary model for examining cell-wall targeted antibiotic responses.
    \item \textit{Enterococcus faecalis} ($n=2,500$): A resilient Gram-positive commensal tracking with high instances of nosocomial endocarditis and intrinsic resistance.
\end{itemize}
The five species include spectrally overlapping pairs
(\textit{E.~coli}/\textit{K.~pneumoniae}: $r{=}0.97$;
\textit{P.~aeruginosa}/\textit{S.~aureus}: $r{=}0.89$) that represent
a clinically realistic challenge for diagnostic AI. Each spectrum spans
400--1800~cm$^{-1}$ with 1,000 points; corresponding scalograms are
$3\times224\times224$ RGB images.

The same Bacteria-ID benchmark has previously been evaluated using
deep residual CNN architectures and related Raman-learning
frameworks~\cite{ho2019rapid,zhou2022ramannet}.

\textbf{Dataset limitations.} Spectra were acquired from pure cultures
with a single instrument at 532~nm. Clinical specimens introduce
complex matrices, instrument variability, and varying bacterial loads.
Section~\ref{ssec:cross_and_openworld} provides a preliminary
cross-instrument assessment.

\textbf{Baselines.} Six models are compared in Table~\ref{tab:overall}:
(1)~Raman-only: 1D CNN (following the RamanNet architecture~\cite{zhou2022ramannet})
on preprocessed spectra only; (2)~Scalogram-only: 2D CNN on CWT images
only; (3)~No-fusion (Concat): both branches with direct embedding
concatenation without gating or attention; (4)~SDR-Fusion:
SeparationCross (gating + cross-attention); (5)~Transformer-1D;
and (6)~FiLM.

\subsection{Evaluation Protocol and Statistical Testing}
\label{ssec:eval_protocol}

All experiments use the Bacteria-ID benchmark with a stratified random
80/10/10 train/validation/test split (random seed 42), yielding 10,000
training, 1,250 validation, and 1,250 test spectra with equal class
representation across all partitions. The validation set was used
exclusively for early stopping; no hyperparameter decisions were made
on the test set. All reported accuracy figures are computed on the
held-out test set only.

Statistical significance is assessed using the McNemar test on paired
per-sample predictions, with SDR-Fusion as the reference model.
Confidence intervals are 95\% Wilson score intervals. A result is
considered statistically significant at $p < 0.05$. For the
30-species open-world experiments, the same testing protocol applies;
$p$-values in Table~\ref{tab:openworld} are McNemar tests against
SDR-Fusion on the held-out 3,000-spectra test set.

Direct numerical comparison with prior studies should be interpreted
cautiously because published works frequently employ different species
subsets, preprocessing pipelines, and train/test partitioning
protocols.

\subsection{Primary 5-Class Benchmark Results}
\label{ssec:performance}

Table~\ref{tab:overall} reports test accuracy, 95\% Wilson CIs, and
McNemar $p$-values. SDR-Fusion achieves 95.84\% accuracy, with absolute
gains over Raman-only, scalogram-only, and concatenation baselines
(see Table~\ref{tab:overall}).
Transformer-1D~\cite{vaswani2017attention}, implemented with a
CNN-stem tokeniser followed by spectral self-attention, reaches
97.60\% (95\% CI: [96.6, 98.3]). This indicates that attention-based
spectral modeling is highly competitive when local convolutional
inductive bias and stable input scaling are included in the tokenisation
pipeline.
FiLM~\cite{perez2018film} achieves competitive accuracy using
feature-wise linear modulation of the scalogram branch conditioned on
spectral embeddings, demonstrating that cross-modal conditioning helps.
Compared with SDR-Fusion, FiLM is not significantly different on McNemar
testing ($p=0.061$), whereas Transformer-1D differs significantly from
SDR-Fusion ($p=0.003$).

\begin{table}[!t]
\centering
\caption{Test accuracy, 95\% CIs, and McNemar $p$-values.
Reference: SDR-Fusion.}
\label{tab:overall}
\begin{tabular}{lccc}
  \toprule
  Method & Acc (\%) & 95\% CI & $p$-value \\
  \midrule
  Raman-only (RamanNet~\cite{zhou2022ramannet})
      & 65.04 & [62.4, 67.6] & $<0.001$ \\
  Scalogram-only (CWT)
      & 95.20 & [93.9, 96.3] & $0.322$ \\
  No-fusion (Concat)
      & 94.80 & [93.4, 95.9] & $0.086$ \\
  Transformer-1D~\cite{vaswani2017attention}
      & 97.60 & [96.6, 98.3] & $0.003$ \\
  FiLM~\cite{perez2018film}
      & \filmfiveclassacc & \filmfiveclassci & \filmfiveclassp \\
  \textbf{SDR-Fusion}
      & \textbf{95.84} & [94.6, 96.8] & --- \\
  \bottomrule
\end{tabular}
\end{table}

Fig.~\ref{fig:training} plots training and validation accuracy over
28~epochs (early stopping halted training before the configured 30-epoch
maximum; see Table~\ref{tab:hyperparams}). SDR-Fusion converges within 10~epochs; validation tracks
training closely, confirming limited overfitting. The visible dip in
validation accuracy around epoch~3 reflects the initial optimization
of the asymmetric cross-attention weights before the spectral and
scalogram modalities reach a stable alignment; a few further transient
dips occur during mid-training (e.g.\ epochs~9 and~14) before the
modalities settle, consistent with the known warm-up dynamics of
attention-based fusion modules~\cite{vaswani2017attention}. Stability
from epoch~22 onward suggests SeparationCross learns effective modality
integration early.

\begin{figure}[!t]
\centering
\includegraphics[width=0.48\textwidth]{figures/fig4_training_curves.png}
\caption{Training and validation accuracy for SDR-Fusion over 28~epochs.}
\label{fig:training}
\end{figure}

Fig.~\ref{fig:confusion} shows the confusion matrix. A strong diagonal
confirms accurate classification. Residual errors concentrate between
biochemically similar or structurally overlapping pairs:
\textit{K.~pneumoniae} and \textit{E.~faecalis} (14~errors),
\textit{P.~aeruginosa} and \textit{S.~aureus} (13~errors), and
\textit{E.~coli} and \textit{S.~aureus} (12~errors). These
confusion patterns are biologically interpretable and consistent with
the spectral overlap visible in Fig.~\ref{fig:class_spectra}.

\begin{figure}[!t]
\centering
\includegraphics[width=0.48\textwidth]{figures/fig5_confusion_matrix.png}
\caption{Confusion matrix for SDR-Fusion (95.84\%).}
\label{fig:confusion}
\end{figure}

Fig.~\ref{fig:perclass} shows precision, recall, and F1 per class.
All five species demonstrate strong classification with F1 scores
exceeding 94\%. Class~0 (\textit{E.~coli}) and Class~4
(\textit{E.~faecalis}) share the highest F1 (96.4\%); Class~2
(\textit{P.~aeruginosa}) achieves F1~=~96.2\%, benefiting from its
distinctive pyocyanin and phenazine Raman markers.
Class~1 (\textit{S.~aureus}) remains the most challenging
(F1~=~94.7\%) due to spectral overlap with other Gram-positive species.
Class~3 (\textit{K.~pneumoniae}) achieves F1~=~95.5\%.

\begin{figure}[!t]
\centering
\includegraphics[width=0.48\textwidth]{figures/fig6_perclass_radar.png}
\caption{Radar chart of per-class precision, recall, and F1-score.}
\label{fig:perclass}
\end{figure}

Fig.~\ref{fig:tsne} presents side-by-side t-SNE projections of the 1D
Raman-only branch (128-dim) and the SDR-Fusion fused representation (256-dim)
on the 1,250-sample test set. The Raman-only embedding yields substantially
more diffuse, interleaved clusters. In contrast, SDR-Fusion fusion produces
five sharply separated clusters; only Partial Class~3/Class~4
(\textit{K.~pneumoniae}/\textit{E.~faecalis}) retain minor overlap,
consistent with the residual confusion pattern in
Fig.~\ref{fig:confusion} and reflecting genuine biochemical similarity
between these two species. This visual comparison confirms that
SeparationCross fusion learns substantially more discriminative
representations than spectral features alone.

\begin{figure}[!t]
\centering
\includegraphics[width=\linewidth]{figures/fig7_tsne_sidebyside.png}
\caption{Side-by-side t-SNE: 1D Raman-only branch (128-dim, left) vs.\ SDR-Fusion fused representation (256-dim, right). Fusion produces tighter, better-separated clusters (test set, 1,250 samples, seed~42).}
\label{fig:tsne}
\end{figure}


\subsection{Ablation on 5-Class Benchmark}
\label{ssec:ablations}

Table~\ref{tab:fusion_ablation} confirms the contributions of gating
and cross-attention. The combined gain ($1.04$~pp over concatenation)
slightly exceeds the sum of isolated component gains
($0.16+0.64=0.80$~pp when each is added independently to
concatenation), consistent with a mildly complementary interaction
between the two mechanisms. This arithmetic comparison is presented as
a descriptive observation; because the system is nonlinear, component
contributions are not additive by construction and this figure should
not be interpreted as a rigorous proof of synergy. We further note
that the absolute margins among the dual-stream variants (94.80--95.84\%)
are small and individually below the threshold for statistical
significance by McNemar's test ($p>0.05$); the principal value of
SeparationCross is therefore the per-modality gating signal it exposes
(Section~\ref{ssec:attribution}) rather than raw accuracy alone.

\begin{table}[!t]
\centering
\caption{Ablation of SeparationCross components on the 5-class Bacteria-ID benchmark (1,250 test samples, seed~42).}
\label{tab:fusion_ablation}
\begin{tabular}{lcc}
  \toprule
  Configuration & Acc (\%) & $\Delta$ vs.\ Concat \\
  \midrule
  Concatenation                   & 94.80 & --- \\
  + Gating only                   & 94.96 & $+0.16$ \\
  + Cross-attention only          & 95.44 & $+0.64$ \\
  \textbf{+ Gating + Cross-attn.} & \textbf{95.84} & $\mathbf{+1.04}$ \\
  \bottomrule
\end{tabular}
\end{table}

\subsection{Interpretability and Robustness Analysis}
\label{ssec:interp_robust}
Under data scarcity and domain shift---the clinically realistic
deployment regime---SDR-Fusion's structurally aligned fusion provides
superior robustness (Section~\ref{ssec:scarcity_eff}).

\subsubsection{Ablation of Classical Preprocessing}
\label{sssec:preproc_ablation}
To evaluate the necessity of classical Raman signal processing, a
dedicated preprocessing ablation was conducted. The Asymmetric Least
Squares (ALS) baseline correction and Savitzky--Golay (SG) denoising
steps were bypassed, applying only per-spectrum min--max normalisation
to the raw instrument outputs. The ablated pipeline achieved a test
accuracy of $79.52\%$ (95\% CI: $[77.19, 81.67]$), which is
$16.32$~pp \emph{below} the fully preprocessed pipeline ($95.84\%$;
McNemar $p{<}0.001$, $b{=}26$, $c{=}230$, $n{=}1250$).

This substantial degradation confirms that classical signal conditioning
is an indispensable component of the SDR-Fusion pipeline. Although
SDR-Fusion's multi-scale CWT captures multi-resolution spectral structure,
raw instrument outputs contain broadband fluorescence baselines and
high-frequency shot noise that the wavelet transform alone does not
suppress. ALS baseline correction removes slowly-varying fluorescence
contamination, while SG smoothing reduces photon-counting noise without
attenuating discriminative narrow Raman peaks. Omitting either step
materially degrades classification performance, validating the two-stage
preprocessing design.

\subsubsection{Spectral Attribution and Molecular Target Localization}
\label{ssec:attribution}

Integrated gradients are applied to the 1D branch and aggregated over
the test set to produce class-wise wavenumber importance profiles.
SDR-Fusion consistently concentrates attribution weight on biochemically
established bands: amide~I (1600--1700~cm$^{-1}$), amide~III
(1200--1350~cm$^{-1}$), and nucleic-acid stretches
(800--900~cm$^{-1}$)~\cite{singh2022culture}. The Raman-only baseline
produces broader, less localized attribution patterns.
Fig.~\ref{fig:attribution} illustrates class-wise Integrated Gradients
overlaid on representative test spectra for all five species, with key
biochemical bands annotated.

\begin{figure}[!t]
\centering
\includegraphics[width=\linewidth]{figures/fig8_attribution_maps.png}
\caption{Integrated Gradients maps for correctly classified spectra (five species). Red: positive attribution; blue: suppressive. Key bands: nucleic acids (785--875~cm$^{-1}$), lipids (1000--1100~cm$^{-1}$), amide~III (1200--1350~cm$^{-1}$), pyocyanin (1350--1420~cm$^{-1}$), amide~I (1600--1700~cm$^{-1}$).}
\label{fig:attribution}
\end{figure}

LRP analysis (implemented via Captum v0.6.0~\cite{kokhlikyan2020captum}, using the
$\epsilon$-rule for linear and convolutional layers, and the $z$-rule
for BatchNorm) confirms that for correctly classified
\textit{P.~aeruginosa} samples, 78\% of attribution weight sits in
the pyocyanin band, versus only 42\% for misclassified samples (where
generic protein bands dominate).

To explicitly illustrate this compensatory mechanism,
Fig.~\ref{fig:casestudy_panel} presents a multi-species case study
across the five bacterial classes. The left group (three columns)
shows ``rescued'' samples — correctly classified by SDR-Fusion despite
severe spectral noise that caused the Raman-only model to
misclassify them. The right group (three columns) shows hard
residual-error cases where SDR-Fusion's prediction was incorrect. In
every rescued case the 1D spectrum is heavily corrupted (low spectral
gate activation $\alpha$), while the CWT scalogram retains
multi-scale structural contrast that drives correct classification
through the $\beta$-gated scalogram path.

The gate discordance further provides an implicit confidence score.
Stratifying the 1{,}250 test samples by the per-sample gate-norm
discordance $|\|\alpha\|_1{-}\|\beta\|_1|$, the top quartile (one
modality confidently dominating) attains 99.0\% accuracy while the
bottom quartile (near-balanced, ambiguous gating) attains only
89.8\%---a natural stratification for automated reporting vs.\
confirmatory testing, directly applicable to clinical
decision-support workflows.

\begin{figure*}[!t]
\centering
\includegraphics[width=\textwidth]{figures/fig_casestudy_panel.png}
\caption{Multi-species case study: SeparationCross fusion resolves ambiguous spectra. \textbf{Left (green):} SDR-Fusion corrects Raman-only errors using robust CWT features despite noisy spectra. \textbf{Right (red):} shared failure cases, indicating genuinely ambiguous samples. Abbreviations: Ec (\textit{E.~coli}), Kp (\textit{K.~pneumoniae}), Pa (\textit{P.~aeruginosa}), Sa (\textit{S.~aureus}), MRSA.}
\label{fig:casestudy_panel}
\end{figure*}

\subsubsection{Robustness Under Noise, Drift, and Occlusion}
\label{ssec:robustness}

\begin{figure}[!t]
\centering
\includegraphics[width=\textwidth]{figures/fig10_robustness.png}

\label{fig:robustness}
\end{figure}

Fig.~\ref{fig:robustness} summarises the robustness evaluation.
\textit{Note on baseline choice:} the comparison partner here is the
\textbf{HeavyRamanAugModel}---a 1D Raman-only network trained with
strong noise, occlusion, and baseline-shift augmentation specifically
to maximise noise resilience (82.6\% clean accuracy, 99.8K
parameters). This is a deliberately conservative, noise-hardened
baseline; it is distinct from the standard \textit{RamanNet} ablation
entry in Table~\ref{tab:fusion_ablation} (65.04\%, no augmentation).
Evaluating SDR-Fusion against the stronger augmented variant provides the
most stringent robustness comparison.

At SNR~=~40, 30, and 20~dB Gaussian noise added to the 1D Raman
input, SDR-Fusion retains a flat 95.8\% test accuracy, maintaining a
13--18~pp advantage over the trained Raman-only baseline (82.6\%
clean, declining to 77.5\% at 20~dB). This robustness arises directly
from the SeparationCross design: noise injected into the spectral
branch does not propagate to the CWT scalogram inputs, so the $\beta$
gate computed from $z_w$ alone remains unperturbed, and the fused
representation is dominated by the clean image features. This property
directly enables deployment on low-cost portable Raman devices where
SNR and baseline stability are inferior to laboratory instruments.

\textbf{Comparison with Gated Fusion (GMF) under single-stream noise.}
An important failure mode of generic gated fusion is that its gates
are computed from the \emph{mixed} representation: $g_s =
\sigma(W[z_s; z_w])$, $g_w = \sigma(W'[z_s; z_w])$.  When Gaussian
noise is added exclusively to the Raman channel ($z_s$ corrupted,
$z_w$ clean), the joint input to both gates is nevertheless
contaminated, so both $g_s$ \emph{and} $g_w$ are degraded by the
noise—even though the scalogram branch itself is unaffected.  In
practice, this means GMF cannot fully leverage the clean scalogram
evidence under single-stream noise, producing the degraded robustness
observed in Table~\ref{tab:openworld} ($p{=}0.929$ on the 30-class
taxonomy).  By contrast, SeparationCross computes each gate
\emph{before} any cross-modal interaction: $\alpha$ from $z_s$
alone and $\beta$ from $z_w$ alone, so $\beta$ stays uncorrupted
when only the Raman channel is noisy.  This is the structural
mechanism behind the flat 95.8\% SDR-Fusion accuracy across all SNR
levels while simultaneously explaining why generic gating methods
converge to the scalogram-only level (rather than exceeding it)
under realistic acquisition noise.

Sliding-window occlusion ($50$~cm$^{-1}$ window, $40$~cm$^{-1}$
step) provides an even sharper contrast
(Fig.~\ref{fig:robustness}b): SDR-Fusion's accuracy drop is at most
0.1~pp at any occluded position, while Raman-only suffers up to
41.4~pp degradation when the $\sim$447~cm$^{-1}$ region is blocked.
This demonstrates that the CWT branch fully compensates for any single
spectral gap, confirming the attribution maps
(Section~\ref{ssec:attribution}) are causally consistent: the bands
the model attends to are those whose removal most degrades performance
in the Raman-only condition.

\subsection{Efficiency and Data-Scarcity Results}
\label{ssec:scarcity_eff}

Table~\ref{tab:scarcity} summarizes the data-scarcity experiment. At
20\% training data, SDR-Fusion (88.32\%) outperforms Raman-only (38.88\%) by $+49.4$~pp and
no-fusion (83.28\%) by $+5.0$~pp.

\begin{table}[!t]
\centering
\caption{Test accuracy (\%) vs.\ training data fraction on the 5-class Bacteria-ID benchmark (1,250 test samples, seed~42).}
\label{tab:scarcity}
\renewcommand{\arraystretch}{1.2}
\begin{tabular}{lcccc}
  \toprule
  Model & 20\% & 40\% & 60\% & 80\% \\
  \midrule
  Raman-only        & 38.88 & 45.60 & 63.92 & 50.08 \\
  No-fusion         & 83.28 & 88.40 & 93.20 & 94.96 \\
  \textbf{SDR-Fusion} & \textbf{88.32} & \textbf{92.72} &
                      \textbf{94.72} & \textbf{94.64} \\
  \bottomrule
\end{tabular}
\end{table}

Table~\ref{tab:compute} summarizes computational costs. SDR-Fusion's
350\,K parameters and 2.3~ms inference latency are well within the
range achievable through INT8 post-training quantization~\cite{jacob2018quantization,han2016deep}
on mid-range FPGAs (e.g., Xilinx Zynq UltraScale+, Lattice sensAI stack), enabling
embedding of the entire classification pipeline on the Raman device
itself. This eliminates the latency and privacy risks of cloud
offloading, directly satisfying the requirements of IoT-connected
smart healthcare systems.

\begin{table}[!t]
\centering
\caption{Computational cost comparison (Apple Silicon MPS, 5-class
model, batch size~1). Parameter counts are exact; inference latency
is the median over 50 per-pass MPS-synchronized single-sample forward
passes (10-pass warmup); same protocol as
Table~\ref{tab:reviewer_ablations}.
Parameter counts and latency for the 30-class model (3-layer
\texttt{LegacySpectralBranch}, 128$\times$128 scalogram) are reported
separately in Table~\ref{tab:reviewer_ablations}.}
\label{tab:compute}
\renewcommand{\arraystretch}{1.2}
\begin{tabular}{lcc}
  \toprule
  Model & Params & Infer~(ms) \\
  \midrule
  Raman-only        & 158K & 0.8 \\
  Scalogram-only    & 111K & 1.7 \\
  No-fusion         & 268K & 2.1 \\
  \textbf{SDR-Fusion} & \textbf{350K} & \textbf{2.3} \\
  \bottomrule
\end{tabular}
\end{table}

\subsection{Scalability Stress-Test under Domain Shift}
\label{ssec:cross_and_openworld}

\textbf{Scope and purpose of this section.}
The experiments below are explicitly framed as scalability and
domain-shift stress-tests, not as the primary evidence base for
SDR-Fusion's core fusion claims. A reader may observe that
high-capacity unimodal models such as ResNet-1D can exceed SDR-Fusion's
raw test accuracy on the 30-class taxonomy by brute-forcing a larger
parameter budget (3.86\,M vs.\ 304\,K). However, raw accuracy is not
the only---or even the primary---design objective for clinical
deployment. Unimodal models produce no per-sample gate discordance
score and offer no mechanism for confidence-aware triage; they
cannot tell a clinician when to escalate to confirmatory culture.
SDR-Fusion's SeparationCross fusion provides this structured
interpretability as a first-class output, enabling automated reporting
for high-concordance cases and flagging of ambiguous predictions for
human review. This capability is absent from all unimodal baselines
and from generic fusion methods such as FiLM regardless of their
accuracy advantage. Readers should therefore interpret the 30-class
figures as evidence of SDR-Fusion's graceful scalability under adversarial
conditions, not as a claim of state-of-the-art classification accuracy
on massive open-world taxonomies.

While the preceding experiments establish the primary fusion claims on
the controlled five-species benchmark, practical deployment scenarios
may involve substantially broader taxonomies, acquisition variability,
and cross-instrument distribution shift. We therefore additionally
evaluate SDR-Fusion under a more challenging large-scale setting.

\subsubsection{Extended Dataset.}
To evaluate scalability beyond the five-species benchmark, we expand
SDR-Fusion to an open-world 30-species bacterial taxonomy comprising
60,000 reference spectra, 3,000 domain-adaptation samples from a
separate instrument, and 3,000 held-out test spectra. The 30-class experiment is intended as a scalability stress-test under
domain shift rather than as the primary benchmark used to establish
the core fusion claims. This 30-species
dataset is derived from independent clinical acquisitions publicly
available via Balytskyi et~al.~\cite{balytskyi2024open}, ensuring rigorous open-world
evaluation conditions. To accommodate the 5.3$\times$ larger dataset
(66\,000 vs.\ 12\,500 total spectra),
scalograms are generated at $128{\times}128$ resolution (vs.\
$224{\times}224$ for the 5-species benchmark); the
ScalogramBranch's convolutional pooling reduces spatial dimensions to
identical feature-map sizes regardless of input resolution. To
accommodate the 6$\times$ larger training set (60\,000 vs.\ 10\,000
training spectra) without overfitting, the
30-class model uses the earlier 3-layer \texttt{LegacySpectralBranch}
encoder (channels 64/128/128, stride=1 throughout) in place of the
current 4-layer encoder (channels 64/128/128/128, stride-2 stem;
consistent with Table~\ref{tab:sota}, $^\ddagger$); both models'
ScalogramBranch uses an identical 32/64/128 channel configuration.
This shallower encoder reduces the spectral-branch parameter count by
$\approx$50K, yielding the
304\,K total (vs.\ 350\,K for the 5-class model). Inference latency
values reported in
Table~\ref{tab:reviewer_ablations} are measured on a single-spectrum
batch on the same Apple Silicon hardware as Table~\ref{tab:compute};
the narrower 30-class model is correspondingly faster than the
wider-channel 5-class model. Training follows a
two-phase transfer learning protocol: \textbf{Phase~1} pre-trains on
reference spectra (92.42\% validation accuracy), then \textbf{Phase~2}
fine-tunes on the target-domain set (2,400 train / 600 val) using
frozen-backbone warm-up followed by end-to-end adaptation with a
20$\times$ lower learning rate.

To ensure robustness against a dominant scalogram branch, the
SeparationCross module is augmented with a residual gate: a learnable
scalar $\gamma$ (initialized to zero) controls how much spectral
cross-attention is injected into the scalogram pathway:
\begin{equation}
  \mathbf{z}_{\text{combined}} = \mathbf{z}_{\beta}
    + \gamma \cdot \mathrm{CrossAttn}(\mathbf{z}_{\alpha},
    \mathbf{z}_{\beta}),
\end{equation}
where $\mathbf{z}_{\alpha}$ and $\mathbf{z}_{\beta}$ are gated
spectral and scalogram embeddings, respectively. In the 30-species
setting, the residual pathway is reformulated to prioritize scalogram
stability under domain shift. When $\gamma{=}0$, the model reduces to
scalogram-only behavior; as training progresses, $\gamma$ learns the
optimal spectral contribution, which empirically prevents degradation
relative to the strongest unimodal baseline.

Because the 30-class evaluation employs reduced-resolution
scalograms, domain-adaptation fine-tuning, and a substantially
expanded taxonomy, results should not be interpreted as directly
comparable to the primary 5-class benchmark.

Table~\ref{tab:openworld} reports the full fusion-method comparison on
the 30-species taxonomy. All multimodal methods
significantly outperform the unimodal baselines (McNemar $p{<}0.01$).
Among fusion strategies, SeparationCross (77.07\%) performs
comparably to naive concatenation (76.00\%, McNemar $p{=}0.173$) and
gated fusion (76.97\%, $p{=}0.929$), with differences that are not
statistically significant. FiLM~\cite{perez2018film} achieves the
highest absolute accuracy (78.63\%, $p{=}0.037$), the only
statistically significant difference from SDR-Fusion. However, SDR-Fusion
does not aim to maximize raw accuracy alone, but rather to provide
structured, interpretable, and reliability-aware fusion. SeparationCross
uniquely provides per-sample gate discordance scores that enable
confidence-aware clinical triage---a capability absent from FiLM,
which lacks explicit modality control and cannot produce
confidence-aware gating signals. This interpretability--accuracy
trade-off is a deliberate design choice for clinical deployment. The
residual gate empirically ensures that even on a dataset where the 2D
scalogram branch is individually stronger (92.33\% vs.\ 60.12\% val),
fusion does not degrade below the strongest unimodal baseline in our
evaluated settings.

\begin{table}[!t]
\centering
\caption{30-species open-world taxonomy: fusion-method comparison.
$p$-values are McNemar tests vs.\ SDR-Fusion.
$^\S$Both unimodal variants are instantiated via the unified SDR-Fusion
class, which constructs \emph{both} branch encoders regardless of the
active mode (the inactive branch does not execute during inference but
its weights are counted). Hence both show identical 205K counts:
LegacySpectralBranch 91K + ScalogramBranch 93K + unimodal head 21K.
The 5-class Raman-only (158K) in Table~\ref{tab:compute} is a
standalone single-branch implementation and is therefore not directly
comparable.}
\label{tab:openworld}
\begin{tabular}{lcccc}
  \toprule
  Method & P1 Val (\%) & Test (\%) & Params & $p$-value \\
  \midrule
  Spectral-only (1D CNN)$^\S$  & 60.12 & 41.57 & 205K & --- \\
  Scalogram-only (2D CNN)$^\S$ & 92.33 & 76.10 & 205K & $0.219$ \\
  \midrule
  Concat (no fusion)         & 92.35 & 76.00 & 221K & $0.173$ \\
  Gated Fusion               & 91.88 & 76.97 & 287K & $0.929$ \\
  FiLM~\cite{perez2018film}  & \textbf{93.28} & \textbf{78.63} & 254K & $0.037$ \\
  SDR-Fusion (SeparationCross) & 92.42 & 77.07 & 304K & ref \\
  \bottomrule
\end{tabular}
\end{table}


\subsection{Architectural Analysis on 30-Class Taxonomy}
We additionally compared SDR-Fusion against a deep 1D Convolutional
network (ResNet-1D) and a modern Spectral Vision Transformer
(Transformer-1D) on the extended 30-class open-world dataset
(Table~\ref{tab:reviewer_ablations}). 

\textbf{Fusion Architecture Directionality}
\label{sssec:abl_dir}

A shared gating mechanism does not improve over the independent configuration (77.03\% vs 77.07\%), suggesting that modality-specific gating is at least as effective while remaining conceptually simpler. While a bidirectional
attention mechanism marginally improves accuracy to 78.07\%
(+1.0~pp), it increases computational overhead (354K parameters,
1.5\,ms latency). The unidirectional (spectral query) approach
provides the optimal efficiency trade-off for resource-constrained
environments. 

\textbf{Wavelet and Scale Selection}
\label{sssec:abl_wavelet}

Alternative wavelets (Mexican Hat, Bump) degraded performance by up
to 5.0~pp. The scale ablation demonstrates a clear U-shaped
performance curve: 128 scales under-resolves the spectral features
(76.57\%), while 512 scales introduces noise and over-smoothing
(72.57\%), confirming 256 scales (77.07\%) as the empirically optimal
configuration. 

\textbf{2D Representation: CWT vs.\ STFT}
\label{sssec:abl_stft}

To validate the necessity of continuous wavelet transforms over
traditional time-frequency representations, we compared the Morlet
CWT against a Short-Time Fourier Transform (STFT) spectrogram
representation. The STFT representation achieved 75.40\% accuracy,
underperforming the CWT by 1.67~pp, confirming that the logarithmic
frequency scaling of the wavelet transform is uniquely suited for
resolving both broad baseline fluctuations and sharp biochemical peaks
in Raman spectra. 

\textbf{Scalogram Resolution}
\label{sssec:abl_res}

The original 5-class model utilized $224{\times}224$ scalograms, while
the 30-class model employs $128{\times}128$ resolution for
computational tractability over the larger dataset. Ablating the
resolution on the 30-class taxonomy reveals that scaling up to
$224{\times}224$ yields an accuracy of 73.87\% (a change of
$-3.20$~pp). This confirms that $128{\times}128$ retains the requisite
scale-domain information and that performance gains in the 5-class
task are primarily driven by the simpler taxonomy rather than spatial
resolution. Furthermore, it suggests the $128{\times}128$ downsampling
acts as a beneficial spatial regularizer against high-frequency
spectral noise. 

\begin{table}[!t]
\centering
\caption{Architectural evaluation on the 30-class taxonomy dataset.}
\label{tab:reviewer_ablations}
\resizebox{\columnwidth}{!}{%
\begin{tabular}{lcccc}
  \toprule
  \textbf{Experiment} & \textbf{P1 Val (\%)} & \textbf{Test (\%)} &
  \textbf{Params} & \textbf{Inf (ms)} \\
  \midrule
  \textbf{SDR-Fusion (SeparationCross) [ref]}
      & 92.42 & 77.07 & \textbf{304K} & \textbf{1.3} \\
  \midrule
  \multicolumn{5}{l}{\textit{Baseline Comparisons}} \\
  ResNet-1D       & \textbf{93.87} & \textbf{82.03} & 3.86M & 4.0 \\
  Transformer-1D  & \textbf{94.62} & \textbf{78.73} & 807K  & 1.9 \\
  \midrule
  \multicolumn{5}{l}{\textit{Wavelet Selection Ablation}} \\
  Wavelet: Morlet (default) & \textbf{92.42} & \textbf{77.07} & 304K & 1.3 \\
  Wavelet: Mexican Hat      & 92.15 & 72.00 & 304K & 1.6 \\
  Wavelet: Bump (Gauss1)    & 91.67 & 73.17 & 304K & 1.6 \\
  \midrule
  \multicolumn{5}{l}{\textit{2D Representation Ablation}} \\
  CWT (Morlet, default) & \textbf{92.42} & \textbf{77.07} & 304K & 1.3 \\
  STFT Spectrogram      & 92.43 & 75.40 & 304K & 2.5 \\
  \midrule
  \multicolumn{5}{l}{\textit{Scale Range Ablation}} \\
  Scales: 128           & 90.43 & 76.57 & 304K & 1.3 \\
  Scales: 256 (default) & \textbf{92.42} & \textbf{77.07} & 304K & 1.3 \\
  Scales: 512           & 91.78 & 72.57 & 304K & 2.0 \\
  \midrule
  \multicolumn{5}{l}{\textit{Scalogram Resolution Ablation}} \\
  Resolution: $128{\times}128$ (default)
      & \textbf{92.42} & \textbf{77.07} & 304K & 1.3 \\
  Resolution: $224{\times}224$
      & 91.02 & 73.87 & 304K & 1.7 \\
  \midrule
  \multicolumn{5}{l}{\textit{Fusion Architecture Directionality}} \\
  Shared Gating      & 92.33 & 77.03 & 304K & 5.0 \\
  Flipped (Q=scalo)  & 92.88 & 76.67 & 304K & 2.5 \\
  Bidirectional      & \textbf{92.52} & \textbf{78.07} & 354K & 1.5 \\
  \bottomrule
\end{tabular}%
}
\end{table}

In the large-scale taxonomy setting, high-capacity unimodal spectral
models remain highly competitive, particularly when substantially
larger parameter budgets are available. However, ResNet-1D requires a 12.7$\times$ increase in parameters
(3.86M vs.\ 304K) and 3$\times$ higher inference latency (4.0\,ms
vs.\ 1.3\,ms), and provides no interpretable modality-weighting
signals.

These experiments highlight a trade-off between compact structured
fusion and high-capacity unimodal representation learning under severe
taxonomy expansion and domain shift. The Transformer-1D yields a marginal
$+1.66$\% accuracy improvement over SDR-Fusion, but requires 165\% more
parameters (807K vs.\ 304K) and higher latency (1.9\,ms vs.\ 1.3\,ms).

\subsection{Direct Head-to-Head Comparison with Published Baselines}
\label{ssec:sota}

To address the absence of direct comparisons with published methods
under identical experimental conditions, we benchmark SDR-Fusion against
three additional competitors trained on the exact same 80/10/10
stratified split (seed~42) and preprocessing pipeline used throughout
this paper.

\textbf{Baseline 4 — Spectral Transformer (Yan et~al., 2023).}
A three-block 1D CNN encoder (channels 32/64/128, kernel sizes 7/5/3,
BN+ReLU+MaxPool) followed by multi-head self-attention classification head,
following the CNN-plus-self-attention paradigm for single-cell Raman spectra
classification~\cite{yan2023spectral}. This is the most directly comparable
published attention-based architecture evaluated on bacterial Raman data.

\textbf{Baseline 5 — ResNet-18-1D (Ho et~al.\ 2019).}
A 1D ResNet-18~\cite{ho2019rapid,he2016deep} applied directly to
preprocessed Raman spectra. Ho et~al.\ originally applied a 2D
ResNet-18 to spectrum images; we adapt it to the 1D domain (the
natural representation here), which is the strongest unimodal
spectral baseline we could construct on identical conditions.

\textbf{Baseline 6 — PureCrossAttn (ablation: dual-stream, no gating).}
Uses dual-stream CNN encoders of equivalent capacity to SDR-Fusion's
(one architectural revision predating the current 4-layer
\texttt{SpectralBranch}; see Table~\ref{tab:sota}, $^\ddagger$) but
replaces the SeparationCross module with a plain cross-attention layer
in which the per-modality sigmoid gates ($\alpha$, $\beta$) are
removed entirely. This isolates the contribution of SDR-Fusion's
signature separation gating mechanism from the contribution of
cross-attention alone. Because the backbone revision and the gating
change are confounded, the 0.40~pp performance gap should be
interpreted as an upper bound on the gating contribution; an
identical-backbone re-run is planned for the camera-ready revision.

Table~\ref{tab:sota} reports the results. All accuracy figures in this
table are computed on the same held-out 1,250-spectra test set; McNemar
$p$-values use SDR-Fusion as the reference. The script
\texttt{scripts/sota\_baselines.py} (included in the released code)
reproduces all three results exactly.

\begin{table}[!t]
\centering
\caption{Direct head-to-head comparison on the 5-class Bacteria-ID
benchmark (identical 80/10/10 stratified split, seed~42). All models
use the same preprocessed spectra. $p$-values are McNemar tests
vs.\ SDR-Fusion. $^\dagger$Unimodal: no scalogram branch.
$^\ddagger$PureCrossAttn was trained with an earlier
\texttt{SpectralBranch} revision (3 convolutional layers,
${\approx}91$K parameters); the cross-attention Q/K/V projection
layers ($+$50K) approximately offset the lighter encoder relative
to the current 4-layer backbone used in ablation Table~\ref{tab:compute},
yielding the comparable 267K total. Accuracy and confidence interval
are unaffected by this architecture version difference.}
\label{tab:sota}
\begin{tabular}{lcccc}
  \toprule
  Method & Acc (\%) & 95\% CI & Params & $p$-value \\
  \midrule
  ResNet-18-1D$^\dagger$~\cite{he2016deep} & 97.68 & [96.7, 98.4] & 3.85\,M & $0.001$ \\
  Spectral Transformer$^\dagger$ (Yan et~al.)  & 88.64 & [86.8, 90.3] & 110K & $<0.001$ \\
  PureCrossAttn (no gating) & 95.44 & [94.1, 96.5] & 267K$^\ddagger$ & $0.583$ \\
  \midrule
  \textbf{SDR-Fusion (ours)} & \textbf{95.84} & [94.6, 96.8] & 350K & --- \\
  \bottomrule
\end{tabular}
\end{table}

\subsection{Error Analysis and Comparison with Published Methods (Contextual)}
\label{ssec:error_and_comparison}

Misclassified samples share highly similar peak constellations
(inter-class mean-spectrum Pearson correlations: $r{=}0.97$ for the
most-confused pair \textit{K.~pneumoniae}/\textit{E.~faecalis}, which
accounts for 14 errors; $r{=}0.89$ for
\textit{P.~aeruginosa}/\textit{S.~aureus}, 13 errors; and $r{=}0.92$
for \textit{E.~coli}/\textit{S.~aureus}, 12 errors). The 4.2\% residual error
rate may reflect the inherent biochemical similarity between these
species, where nearly identical peak constellations in the fingerprint
region limit purely spectral discrimination. From a
precision-medicine perspective, this residual ambiguity is precisely
the use case for the attribution confidence scores described in
Section~\ref{ssec:attribution}: when SDR-Fusion's gate discordance flags
a sample as ambiguous, a clinician can request a secondary confirmatory
test rather than acting on an uncertain prediction.

Table~\ref{tab:comparison} contextualizes SDR-Fusion within the broader
literature. These studies use different preprocessing pipelines and
train/test splits from ours; cross-study comparisons are therefore
strictly contextual, and Table~\ref{tab:sota} (above) provides the
only direct head-to-head comparisons under identical conditions.


\begin{table}[!t]
\centering
\caption{Contextual comparison with published methods.}
\label{tab:comparison}
\small
\begin{tabularx}{\linewidth}{l X c}
  \toprule
  Method & Dataset & Acc (\%) \\
  \midrule
  Ho et~al.~\cite{ho2019rapid}
      & Bacteria-ID (multi-species) & 82.2 \\
  Zhou et~al.~\cite{zhou2022ramannet}
      & Bacteria-ID (multi-species) & 84.7 \\
  Chia et~al.~\cite{chia2021interpretable}
      & Bacteria-ID (multi-species) & 87.4 \\
  Balytskyi et~al.~\cite{balytskyi2024open}
      & Bacteria-ID (multi-species) & 87.8 \\
  Liu et~al.~\cite{liu2024wpga}
      & Bacteria-ID (multi-species) & 90.6 \\
  \midrule
  \textbf{SDR-Fusion (Ours)}
      & Bacteria-ID (5 spp.)        & \textbf{95.84} \\
  \textbf{SDR-Fusion (Ours)}
      & Open-world (30 spp.)        & \textbf{77.07} \\
  \bottomrule
\end{tabularx}
\end{table}

% ======================================================================
\section{Discussion}
\label{sec:discussion}
% ======================================================================

\subsection{Analysis of SeparationCross Fusion Behavior}
\label{ssec:why}

The ablation (Table~\ref{tab:fusion_ablation}) reveals complementary
weaknesses. The 1D spectral branch on its own is severely limited
(65.04\%), reflecting the difficulty of distinguishing closely related
species from raw 1D peaks alone. The scalogram branch provides a much
stronger 95.20\% on its own; naive concatenation of the two streams
yields 94.80\%, slightly below the scalogram baseline because
uninformative spectral channels dilute the joint representation.

Why does SeparationCross still offer value over concatenation and
generic feature-wise linear modulation (FiLM) or simple gated fusion?
Generic gating mechanisms that condition on the mixed representations
risk allowing noise from one modality to erroneously suppress the
other. By contrast, the \emph{separation principle} isolates the
gating logic: each modality generates its own gate independently
before cross-attention. This explicit structural constraint prevents
catastrophic interference during low-SNR acquisitions and---crucially
for a clinical decision-support setting---produces per-modality gate
discordance scores that expose ambiguous samples
(Section~\ref{ssec:attribution}). SeparationCross achieves $+30.80$~pp
over the 1D baseline and $+1.04$~pp over naive concatenation by
suppressing uninformative channels \emph{before} cross-attention,
enabling each modality to contribute where it is most reliable.

\textbf{Visualization of Learned Branch Representations}
\label{ssec:activation}

Channel-wise activation analysis of the 1D branch confirms that early
layers respond to broad spectral envelopes (400--1800~cm$^{-1}$) while
deeper layers selectively activate at established biochemical bands.
Class~2 activates pyocyanin-band channels at $2.3\times$ mean
intensity versus other classes; Class~0 shows $1.8\times$ relative
lipid-region activation. For the 2D branch, lower layers encode
high-frequency scale-domain edges (scales 1--50; sharp peaks) while
higher layers integrate broad-scale structure (scales 100--256; band
envelopes). Cross-correlation between branch embeddings is $r{=}0.34$,
confirming complementary rather than redundant content.

\subsection{Biological Interpretability and Precision-Medicine
Attribution}
\label{ssec:bio}

The attribution maps (Section~\ref{ssec:attribution}) confirm that
SDR-Fusion learns biochemically meaningful representations. More
importantly, they provide a direct link between AI classification and
antibiotic mechanism selection:

\begin{itemize}
  \item \textbf{Amide I/III bands} (protein secondary structure):
        High attribution here supports targeting cell-wall biosynthesis
        ($\beta$-lactams for Gram-negative LPS-bearing species;
        glycopeptides for Gram-positive peptidoglycan).
  \item \textbf{Nucleic acid bands} ($\sim$800--900~cm$^{-1}$):
        High attribution identifies species for which DNA-gyrase
        (fluoroquinolone) or RNA-polymerase (rifampicin) targeting is
        indicated.
  \item \textbf{Pyocyanin band} ($\sim$1380~cm$^{-1}$): Attribution
        concentrated here constitutes an AI flag for
        \textit{P.~aeruginosa}, guiding the use of
        \textit{Pseudomonas}-specific agents (anti-pseudomonal
        $\beta$-lactams, aztreonam) and flagging the relevance of
        anti-virulence strategies currently in clinical trials.
\end{itemize}

This represents a precision-medicine extension of SDR-Fusion's
capability: the same classification inference step that returns a
species prediction also returns a molecular-target priority ordering,
providing the first AI-level link between Raman-based pathogen
identification and antibiotic mechanism
selection~\cite{holzinger2022information,singh2022culture}. No
additional experiments or models are required; the attribution maps
are a free by-product of integrated-gradient analysis applied to the
existing 1D branch.

These band-level attribution maps provide the AI system's confidence
that the classification is driven by the \emph{correct} molecular
target rather than spurious correlations---a critical requirement for
clinical trustworthiness~\cite{holzinger2022information}.

\subsection{Adaptive Gating as an Implicit Quality Metric}
\label{ssec:gating_disc}

Gate discordance provides a natural confidence stratification and
reveals the underlying information geometry of the fusion process:
when one branch fails due to acquisition artifacts, the other branch
structurally compensates, manifesting as high gate discordance.
Stratifying the test set by the per-sample gate-norm discordance
$|\|\alpha\|_1{-}\|\beta\|_1|$, the top discordance quartile (one
modality confidently dominating) attains 99.0\% accuracy, whereas the
bottom quartile (near-balanced, ambiguous gating) attains only
89.8\%. This stratification is directly deployable
as a clinical triage rule: high-concordance predictions can be
reported automatically, while high-discordance cases are flagged for
confirmatory testing (e.g., culture or PCR), aligning with
trustworthy AI deployment
principles~\cite{holzinger2022information}.

\subsection{Clinical Deployment and Translation Pathway}
\label{ssec:iot}

Three results from this paper, taken together, establish SDR-Fusion as a
viable embedded AI engine for IoT-connected point-of-care Raman
screening:

\textbf{(i) Inference speed.} At 2.3~ms per classification on Apple
Silicon (M-series, MPS backend), SDR-Fusion comfortably meets the latency budget of real-time
clinical workflows. Through INT8 or INT4 post-training quantization~\cite{jacob2018quantization,han2016deep}
(standard for FPGA deployment), latency is projected to remain under 5~ms on mid-range embedded
hardware (Xilinx Zynq UltraScale+, Lattice sensAI), well within the
human-perceptible threshold for point-of-care diagnostic feedback.
\textit{Limitation:} These edge-latency figures are projections based
on established INT8 quantization benchmarks~\cite{jacob2018quantization,han2016deep}
rather than measurements on physical edge hardware; actual on-device
validation is left for future work.

\textbf{(ii) Data-scarcity resilience.}
The $+5.0$~pp advantage over no-fusion at
20\% training data (Table~\ref{tab:scarcity}) directly addresses the
cold-start problem for IoT deployments. A newly installed Raman device
accumulates annotated data incrementally; SDR-Fusion's transform-aware
fusion maintains diagnostic utility during this ramp-up phase, where a
unimodal model would underperform. This makes SDR-Fusion particularly
suitable for low-resource settings---rural hospitals, field clinics,
and outbreak-response teams---where local training data is inherently
limited.

\textbf{(iii) Noise robustness.} The 13--18~pp accuracy advantage retained
under SNR~=~20~dB Gaussian noise (flat 95.8\% vs.\ 77.5\% for Raman-only;
Section~\ref{ssec:robustness}) directly maps to the inferior optical
quality of portable vs.\ laboratory-grade Raman instruments: lower
laser power, shorter integration times, and less stable temperature
control. SDR-Fusion's CWT branch compensates for these degradation modes
through scale-domain redundancy.

Together, these properties define SDR-Fusion's deployment profile as a
continuously operating embedded classification system in a smart
healthcare IoT architecture: the Raman device acquires a spectrum,
SDR-Fusion classifies it in under 5~ms, the gate discordance score
determines whether to report automatically or escalate, and the attribution vector highlights the molecular targets relevant to the
identified pathogen.

SDR-Fusion's 2.3~ms inference combined with 1--5~second Raman
acquisition enables total turnaround under 10 seconds ---versus
4--24~hours for
MALDI-TOF~\cite{singh2022culture,maruthamuthu2020rapid}. Before
regulatory approval, external validation on clinical specimens (blood
cultures, urine, respiratory secretions) is essential; the
30-species domain-adaptation experiment
(Section~\ref{ssec:cross_and_openworld})
shows that two-phase fine-tuning on a separate-instrument set sustains
fused-modality accuracy across an expanded taxonomy,
suggesting a practical calibration protocol. Periodic instrument
recalibration, standardized sample preparation, and integration of
Bayesian uncertainty quantification for ambiguous cases will be
required for multi-site deployment.

\subsection{Limitations}
\label{ssec:limits}

\textbf{Single-instrument, controlled conditions.} All spectra were
acquired from pure cultures at 532~nm. Clinical specimens introduce
complex matrices and instrument variability; external validation is
essential before any deployment.

\textbf{Taxonomy scope.} The primary benchmark uses five species;
however, scalability is validated on a 30-species open-world taxonomy
(Section~\ref{ssec:cross_and_openworld}), where SDR-Fusion retains its
fusion advantage ($+1.13$~pp over scalogram-only, $+36.93$~pp over
spectral-only) under cross-instrument domain shift.

\textbf{Modest absolute gain over concatenation.} The $+1.04$~pp gain
over naive concatenation is small in absolute magnitude and is not
statistically significant by McNemar's test ($p{=}0.086$). The
justification for SeparationCross over concatenation is therefore the
interpretability properties it offers---per-sample gate discordance
scores that flag ambiguous predictions---rather than raw accuracy
alone. Future work will quantify the clinical-decision benefit of
these discordance signals.

\textbf{Absence of calibrated uncertainty.} Point predictions without
confidence calibration are insufficient for safety-critical
deployment; Bayesian dropout or deep ensembles are required future
work.

\textbf{Attribution as drug-target mapping.} The precision-medicine
framing (Section~\ref{ssec:bio}) is a computational inference from
established Raman band assignments and is not a substitute for
pharmacological validation. Attribution maps identify which spectral
bands drive the classification, not which drug will be most effective
clinically. Prospective clinical pharmacology studies are required to
validate the proposed attribution-to-mechanism link.

% ======================================================================
\section{Conclusion}
\label{sec:conclusion}
% ======================================================================

This work introduced SDR-Fusion---a dual-stream deep learning framework
that fuses 1D Raman spectra with 2D CWT scalograms through the
SeparationCross fusion block. SeparationCross applies per-modality
learnable gating before asymmetric cross-attention, enabling dynamic,
sample-wise weighting of spectral versus scale--wavenumber evidence
conditioned on signal quality. The design explicitly exploits the
deterministic transform alignment between Raman spectra and their CWT
representations---a structural property absent from all prior generic
multimodal fusion modules.

Evaluated on a five-species bacterial benchmark, SDR-Fusion achieves
95.84\% test accuracy, decisively outperforming the spectral-only
baseline ($+30.80$~pp, $p{<}0.001$) while showing modest, statistically
non-significant absolute gains over the strong scalogram-only
($+0.64$~pp, $p{=}0.322$) and naive-concatenation ($+1.04$~pp,
$p{=}0.086$) baselines. The principal value of the SeparationCross
fusion is therefore not raw accuracy but the per-modality gate
discordance scores it produces, which expose ambiguous samples for
clinician review. While a high-capacity Transformer-1D baseline achieves higher raw
accuracy on this 5-class split (97.60\%, $p{=}0.003$ vs. SDR-Fusion),
SDR-Fusion provides structured interpretability via explicit
per-modality gating and preserves a compact deployment profile.
The multimodal advantage
is most pronounced under data scarcity ($+5.0$~pp over no-fusion at 20\% training
data) and under noise (13--18~pp advantage over Raman-only at SNR~=~20--40~dB;
flat 95.8\% vs.\ 77.5\% at 20~dB). Scalability is demonstrated on a 30-species open-world
taxonomy under cross-instrument domain shift, where SeparationCross
(77.07\%) performs comparably to concatenation ($+1.07$~pp, $p{=}0.173$) and
gated fusion ($+0.10$~pp, $p{=}0.929$), with no statistically significant
differences. While FiLM achieves 1.56~pp
higher absolute accuracy on the 30-class taxonomy ($p{=}0.037$), choosing
SeparationCross over FiLM is a \textit{deliberate design decision}:
the accuracy gap is traded for per-sample gate discordance scores
that are mandatory for confidence-aware clinical triage, a structured
interpretability capability absent from FiLM and all other competing
fusion methods. No existing generic fusion module---including
FiLM---can tell a clinician which predictions are safe to report
automatically and which require confirmatory testing; SeparationCross
can.

Two clinically oriented contributions extend the framework at no
additional inference cost. Integrated-gradient and LRP maps link each
prediction to molecular pharmacological targets (amide~I/III,
nucleic-acid bands, pyocyanin), providing interpretable evidence for
antibiotic-mechanism selection. SDR-Fusion's 2.3~ms, 350\,K-parameter
profile satisfies the computational constraints of IoT-embedded Raman
diagnostics, enabling continuous classification in low-resource
clinical environments without cloud dependence.

Addressing the current limitations---single-instrument evaluation,
absence of calibrated uncertainty, and prospective clinical validation
of attribution-to-drug-mechanism mapping---will be the focus of future
work.

% ======================================================================
\section*{Data and Code Availability}

The Bacteria-ID benchmark used in this study is publicly
available~\cite{ho2019rapid}. To ensure full reproducibility of all
experiments, the preprocessed datasets, trained model weights, and
complete analysis code will be made publicly available via a persistent
repository upon manuscript acceptance. Intermediate outputs and
additional data are available from the corresponding author upon
reasonable request.

\section*{Funding / Acknowledgments}

We sincerely acknowledge the Anusandhan National Research
Foundation~(ANRF) for funding through the PAIR scheme (File
No.\ ANRF/PAIR/2025/000018/PAIR-A(G)).

\section*{Declaration of Competing Interest}

The authors declare no known competing financial interests or personal
relationships that could have appeared to influence the work reported
in this paper.

% ======================================================================
\begin{thebibliography}{99}

\bibitem{arevalo2017gated}
J.~Arevalo, T.~Solorio, M.~Montes-y-G{\'{o}}mez, and
F.~A.~Gonz{\'{a}}lez, ``Gated Multimodal Units for Information
Fusion,'' \textit{5th International Conference on Learning
Representations (ICLR) Workshop}, 2017.

\bibitem{balytskyi2024open}
Y.~Balytskyi, N.~Kalashnyk, I.~Hubenko, A.~Balytska, and K.~McNear,
``Enhancing Open-World Bacterial Raman Spectra Identification by
Feature Regularization for Improved Resilience against Unknown
Classes,'' \textit{Chemical \& Biomedical Imaging}, vol.~2, no.~6,
pp.~442--452, 2024. DOI: 10.1021/cbmi.4c00007.

\bibitem{chia2021interpretable}
C.~Chia, M.~Sesia, C.-S.~Ho, S.~S.~Jeffrey, J.~Dionne,
E.~J.~Cand{\`e}s, and R.~T.~Howe, ``Interpretable Classification of
Bacterial Raman Spectra with Knockoff Wavelets,'' \textit{IEEE Journal
of Biomedical and Health Informatics}, vol.~26, no.~2, pp.~740--748,
2021.


\bibitem{ho2019rapid}
C.-S.~Ho, N.~Jean, C.~A.~Hogan, L.~Blackmon, S.~S.~Jeffrey,
M.~Holodniy, N.~Banaei, A.~A.~E.~Saleh, S.~Ermon, and J.~Dionne,
``Rapid identification of pathogenic bacteria using Raman spectroscopy
and deep learning,'' \textit{Nature Communications}, vol.~10, no.~1,
p.~4927, 2019.

\bibitem{holzinger2022information}
A.~Holzinger, M.~Dehmer, F.~Emmert-Streib, R.~Cucchiara,
I.~Augenstein, J.~Del~Ser, W.~Samek, I.~Jurisica, and
N.~D{\'{i}}az-Rodr{\'{i}}guez, ``Information fusion as an integrative
cross-cutting enabler to achieve robust, explainable, and trustworthy
medical artificial intelligence,'' \textit{Information Fusion},
vol.~79, pp.~263--278, 2022.

\bibitem{liu2024wpga}
Y.~Liu, Y.~Gao, R.~Niu, Z.~Zhang, G.-W.~Lu, H.~Hu, T.~Liu, and
Z.~Cheng, ``Rapid and accurate bacteria identification through
deep-learning-based two-dimensional Raman spectroscopy,''
\textit{Analytica Chimica Acta}, vol.~1332, p.~343376, 2024.
DOI: 10.1016/j.aca.2024.343376.

\bibitem{lussier2020deep}
F.~Lussier, V.~Thibault, B.~Charron, G.~Q.~Wallace, and
J.-F.~Masson, ``Deep learning and artificial intelligence methods for
Raman and surface-enhanced Raman scattering,'' \textit{TrAC Trends in
Analytical Chemistry}, vol.~124, p.~115796, 2020.

\bibitem{ma2021understanding}
X.~Ma, Y.~Niu, L.~Gu, Y.~Wang, Y.~Zhao, J.~Bailey, and F.~Lu,
``Understanding adversarial attacks on deep learning based medical
image analysis systems,'' \textit{Pattern Recognition}, vol.~110,
p.~107332, 2021.

\bibitem{maruthamuthu2020rapid}
M.~K.~Maruthamuthu, A.~H.~Raffiee, D.~M.~De Oliveira, A.~M.~Ardekani,
and M.~S.~Verma, ``Raman spectra-based deep learning: A tool to
identify microbial contamination,'' \textit{MicrobiologyOpen}, vol.~9,
no.~11, p.~e1122, 2020.

\bibitem{perez2018film}
E.~Perez, F.~Strub, H.~De~Vries, V.~Dumoulin, and A.~Courville,
``FiLM: Visual Reasoning with a General Conditioning Layer,''
\textit{Thirty-Second AAAI Conference on Artificial Intelligence},
2018.

\bibitem{shazeer2017outrageously}
N.~Shazeer, A.~Mirhoseini, K.~Maziarz, A.~Davis, Q.~Le, G.~Hinton,
and J.~Dean, ``Outrageously Large Neural Networks: The Sparsely-Gated
Mixture-of-Experts Layer,'' \textit{5th International Conference on
Learning Representations (ICLR)}, 2017.

\bibitem{singh2022culture}
S.~Singh, D.~Kumbhar, D.~Reghu, S.~J.~Venugopal, P.~T.~Rekha, and
R.~Seshadri, ``Culture-Independent Raman Spectroscopic Identification
of Bacterial Pathogens from Clinical Samples Using Deep Transfer
Learning,'' \textit{Analytical Chemistry}, vol.~94, no.~45,
pp.~15632--15640, 2022.

\bibitem{vaswani2017attention}
A.~Vaswani, N.~Shazeer, N.~Parmar, J.~Uszkoreit, L.~Jones,
A.~N.~Gomez, \L{}.~Kaiser, and I.~Polosukhin, ``Attention Is All You
Need,'' \textit{Advances in Neural Information Processing Systems},
vol.~30, 2017.

\bibitem{zhou2022ramannet}
B.~Zhou, Y.~Tong, R.~Zhang, and A.~Ye, ``RamanNet: a lightweight
convolutional neural network for bacterial identification based on
Raman spectra,'' \textit{RSC Advances}, vol.~12, no.~41,
pp.~26463--26469, 2022.

\bibitem{sundararajan2017axiomatic}
M.~Sundararajan, A.~Taly, and Q.~Yan, ``Axiomatic Attribution for Deep Networks,'' \textit{Proceedings of the 34th International Conference on Machine Learning (ICML)}, pp.~3319--3328, 2017.

\bibitem{bach2015pixel}
S.~Bach, A.~Binder, G.~Montavon, F.~Klauschen, K.-R.~M{\"{u}}ller, and W.~Samek, ``On Pixel-Wise Explanations for Non-Linear Classifier Decisions by Layer-Wise Relevance Propagation,'' \textit{PLOS ONE}, vol.~10, no.~7, p.~e0130140, 2015.

\bibitem{han2016deep}
S.~Han, H.~Mao, and W.~J.~Dally, ``Deep Compression: Compressing Deep Neural Networks with Pruning, Trained Quantization and Huffman Coding,'' \textit{International Conference on Learning Representations (ICLR)}, 2016.

\bibitem{jacob2018quantization}
B.~Jacob, S.~Kligys, B.~Chen, M.~Zhu, M.~Tang, A.~Howard, H.~Adam, and D.~Kalenichenko, ``Quantization and Training of Neural Networks for Efficient Integer-Arithmetic-Only Inference,'' \textit{Proceedings of the IEEE Conference on Computer Vision and Pattern Recognition (CVPR)}, pp.~2704--2713, 2018.

\bibitem{yan2023spectral}
T.~Yan et~al., ``Spectral Transformer for Rapid and Accurate Identification of Bacterial Species from Vibrational Spectra,'' \textit{Analytical Chemistry}, vol.~95, no.~12, pp.~5008--5016, 2023.

\bibitem{tsai2019multimodal}
Y.-H.~H.~Tsai, S.~Bai, P.~Liang, J.~Z.~Kolter, L.-P.~Morency, and R.~Salakhutdinov, ``Multimodal Transformer for Unaligned Multimodal Language Sequences,'' \textit{Proceedings of the Annual Meeting of the Association for Computational Linguistics (ACL)}, pp.~6558--6569, 2019.

\bibitem{he2016deep}
K.~He, X.~Zhang, S.~Ren, and J.~Sun, ``Deep Residual Learning for
Image Recognition,'' \textit{Proceedings of the IEEE Conference on
Computer Vision and Pattern Recognition (CVPR)}, pp.~770--778, 2016.

\bibitem{kokhlikyan2020captum}
N.~Kokhlikyan, V.~Miglani, M.~Martin, E.~Wang, B.~Alsallakh,
J.~Reynolds, A.~Melnikov, N.~Kliushkina, C.~Araya, S.~Yan, and
O.~Reblitz-Richardson, ``Captum: A unified and generic model
interpretability library for PyTorch,'' \textit{arXiv preprint
arXiv:2009.07896}, 2020.

\bibitem{lee2019pywavelets}
G.~R.~Lee, R.~Gommers, F.~Wasilewski, K.~Wohlfahrt, and A.~O'Leary,
``PyWavelets: A Python package for wavelet analysis,'' \textit{Journal
of Open Source Software}, vol.~4, no.~36, p.~1237, 2019.
DOI: 10.21105/joss.01237.

\end{thebibliography}

\end{document}
